% concatenated from main_arxiv.tex
\documentclass{article}
\usepackage[a4paper, left=0.8in, right=0.8in, top=1in, bottom=1in]{geometry}
% The option smallextended is the standard JOTA format.
% The option referee  makes the paper double-spaced.
% The option envcountsect numbers theorems, etc, by section.
% svjour3 is the document class for Springer journals.

%This command right justifies   throughout the paper.
\usepackage{graphicx}
\usepackage{caption}
\captionsetup{width=0.9\textwidth} % Limits caption width
\usepackage{enumerate,comment,mathtools,amsmath,amssymb,algorithm,amsthm}
\usepackage{algpseudocode}
\newtheorem{lemma}{Lemma}[section]
\newtheorem{theorem}[lemma]{Theorem}
\newtheorem{corollary}[lemma]{Corollary}
\newtheorem{definition}{Definition}[section]
\newtheorem{proposition}{Proposition}[section]
\newtheorem{cor}{Corollary}[section]
\newtheorem{remark}{Remark}[section]
\newtheorem{example}{Example}[section]
\newtheorem{setting}{Setting}[section]

\usepackage{xcolor}
\newcommand{\first}[1]{{\color{black}{#1}}}
\newcommand{\second}[1]{{\color{black}{#1}}}
\newcommand{\both}[1]{{\color{black}{#1}}}

\begin{document}

\title{The State-Dependent Riccati Equation in Nonlinear Optimal Control: Analysis and Numerical Approximation}


\author{
    Luca Saluzzi \thanks{Department of Mathematics and Computer Science, University of Ferrara, Italy (\texttt{luca.saluzzi@unife.it})}
}


\maketitle
% -- Abstract
\begin{abstract}
The State-Dependent Riccati Equation (SDRE) approach is extensively utilized in nonlinear optimal control as a reliable framework for designing robust feedback control strategies. This work provides an analysis of the SDRE approach, examining its theoretical foundations, error bounds, and numerical approximation techniques. We explore the relationship between SDRE and the Hamilton-Jacobi-Bellman (HJB) equation, deriving residual-based error estimates to quantify its suboptimality. Additionally, we introduce an optimal semilinear decomposition strategy to minimize the residual. From a computational perspective, we analyze two numerical methods for solving the SDRE: the offline–online approach and the Newton–Kleinman iterative method. Their performance is assessed through a numerical experiment involving the control of a nonlinear reaction-diffusion PDE. Results highlight the trade-offs between computational efficiency and accuracy, indicating better performance of the Newton–Kleinman approach in achieving stable and cost-effective solutions in the reported experiments.
\end{abstract}




\bigskip


\section{Introduction}

    
Optimal control of nonlinear dynamical systems is a fundamental problem in engineering, economics, and applied mathematics. The Hamilton-Jacobi-Bellman equation provides a rigorous framework for computing optimal feedback control laws, whose foundational work dates back to Bellman’s seminal contributions \cite{bellman1957}, establishing dynamic programming as a cornerstone of optimal control theory. However, solving the HJB equation in practical scenarios remains intractable due to its nonlinear and high-dimensional nature. To mitigate this issue, various approximation techniques have been developed, including sparse grids strategies \cite{GK16}, polynomial approximations \cite{kunisch2023learning,kalise2018polynomial}, applications of artificial neural networks \cite{Darbon_Langlois_Meng_2020,Kunisch_Walter_2021,Zhou_2021} and regression-type methods in tensor formats \cite{oster22,dolgov2023data}.
In this context, the State-Dependent Riccati Equation method has emerged as an effective approach, offering a computationally feasible alternative by leveraging local approximations of the value function through a Riccati-based formulation. The SDRE method extends the classical Linear Quadratic Regulator (LQR) framework to nonlinear systems by representing the system dynamics in a state-dependent linearized form. This transformation enables the application of Riccati-based feedback synthesis, preserving the stabilizing properties of LQR while maintaining adaptability to nonlinear dynamics. A key advantage of SDRE is its ability to provide suboptimal yet computationally efficient solutions, making it particularly useful in applications \cite{voos2006nonlinear,vaddi2009numerical,heiland2023low}.
 %Among these, the SDRE approach has gained significant attention for its balance between computational efficiency and solution accuracy. 
In \cite{banks2007}, the authors developed a rigorous formulation of SDRE-based control, demonstrating its effectiveness in stabilizing nonlinear systems under mild regularity assumptions. Moreover, it has been shown that the SDRE feedback induces a stationary Hamiltonian with respect to the control, ensuring that as the state converges to zero, the necessary optimality conditions are asymptotically satisfied at a quadratic rate. For a comprehensive discussion on this topic, we refer to \cite{CIMEN,cloutier1996nonlinear}.
% More recent studies have explored enhancements to the SDRE methodology, including optimization-based formulations to minimize approximation errors \cite{jones2020,dolgov2022optimizing}.
Despite its advantages, the SDRE approach is not without limitations. The choice of the semilinear representation can significantly impact the accuracy of the solution, and existing methods often rely on heuristics for selecting an appropriate decomposition. 
 %Recent research has sought to address this issue by formulating optimal semilinear decompositions that minimize the residual error between the SDRE approximation and the true value function \cite{ran1988parameter}. 
 Furthermore, numerical solvers for the SDRE require efficient methodologies for iteratively solving the associated Riccati equations. An initial attempt to mitigate the computational complexity is presented in \cite{beeler2000feedback}, where the authors introduce a power series expansion for the SDRE solution. This approach involves computing certain terms offline by solving a Riccati equation alongside a sequence of Lyapunov equations. However, for high-dimensional problems, this method may become impractical due to the computational burden and storage requirements associated with high-dimensional matrix equations. To address this challenge, the authors of \cite{alla2023state} proposed an \emph{offline-online} approach based on a first-order approximation, wherein the computational effort is partitioned into an intensive offline phase and a more efficient online stage. Another promising avenue involves solving the SDRE using the Newton-Kleinman method \cite{kleinman1968,saluzzi2025dynamical}, which capitalizes on the iterative nature of the approximation process.
 %Furthermore, numerical solvers for the SDRE require efficient techniques for solving the associated Riccati equations iteratively. A first attempt to reduce the computational complexity can be found in \cite{beeler2000feedback}, where the authors propose a power series expansion for the solution of the SDRE, involving terms that can be computed offline solving a Riccati equation and a series of Lyapunov equations. For high-dimensional problems this can be unfeasible due to the computation and the storage of high dimensional matrix equations. For this reason in \cite{alla2023state} the authors proposed an \emph{offline-online} approach based on a first order approximation, where the computational effort has been decoupled on a heavy offline phase and a cheaper online stage. Another promising direction relies on the resolution via the Newton-Kleinman method \cite{kleinman1968,saluzzi2025dynamical}, taking advantage of the iterative nature of the approximation.
 
 %The Newton-Kleinman method \cite{kleinman1968} and the offline-online strategy \cite{alla2023state,beeler2000feedback} represent two prominent approaches, each with trade-offs in computational efficiency and accuracy.
 
\subsection{Contributions and Outline}
This paper aims to provide a deeper understanding of the SDRE method by analyzing its theoretical properties, deriving error bounds for the approximation, and exploring computational techniques for its implementation. Our contributions include:
\begin{itemize}
\item a derivation of error bounds for the SDRE approximation, quantifying the deviation from the optimal HJB solution;
\item an investigation of optimal semilinear representations that minimize the SDRE residual error,
\item a comparative study of numerical methods for solving the sequence of Riccati equations, highlighting the advantages of iterative methods.
\end{itemize}
The paper is structured as follows: Section 2 formulates the infinite-horizon optimal control problem and introduces key concepts related to the SDRE framework. Section 3 examines error bounds for the SDRE approach, focusing on the residual with respect to the HJB equation and the existence of an optimal semilinear formulation. Section 4 presents numerical techniques for solving the SDRE, including a comparative analysis of two computational methods applied to the control of a nonlinear reaction-diffusion PDE. Finally, Section 5 concludes with a discussion on potential directions for future research.

\begin{comment}
    
Optimal control of nonlinear dynamical systems constitutes a central challenge across diverse fields encompassing engineering disciplines, economic modeling, and applied mathematics. The quest for optimal feedback control laws has been rigorously formalized by the Hamilton-Jacobi-Bellman (HJB) equation. This cornerstone of optimal control theory, originating from Bellman's groundbreaking work on dynamic programming \cite{bellman1957dynamic}, provides a complete theoretical framework.  The HJB equation, derived from Bellman's principle of optimality, dictates the evolution of the optimal value function, which in turn yields the desired optimal control policy through differentiation.  However, despite its theoretical elegance and completeness, the practical applicability of directly solving the HJB equation is severely limited. This intractability arises primarily from the inherent nonlinearity and, critically, the exponential growth in computational complexity with the system's state-space dimension—a phenomenon commonly known as the "curse of dimensionality" \cite{bellman2015dynamic}.  The HJB equation, typically a nonlinear partial differential equation (PDE), becomes exceedingly challenging to solve analytically for even moderately complex nonlinear systems, and numerical discretization further exacerbates the computational burden, rendering exact solutions unattainable in most realistic scenarios.

To circumvent these formidable challenges, a plethora of approximation techniques have been developed to obtain practically usable, albeit often suboptimal, solutions to the optimal control problem.  These approaches can be broadly categorized, and include methods leveraging sparse grid discretizations to mitigate the curse of dimensionality in numerical PDE solvers \cite{GK16sparsegrids}, techniques employing artificial neural networks (ANNs) to approximate the value function and control policy through function approximation and learning algorithms \cite{Darbon_Langlois_Meng_2020ANNhjb,Kunisch_Walter_2021deepHJB,Zhou_2021policyiterationANN}.  For instance, neural network-based approaches can parameterize the value function and control policy and optimize network weights to satisfy HJB equation residuals, offering a powerful tool for high-dimensional problems. Regression-type methods, particularly those employing tensor decompositions, have also emerged as promising avenues for dimensionality reduction and efficient computation. These tensor-based methods, such as those based on tensor train or hierarchical Tucker formats, aim to represent high-dimensional functions and operators in compressed formats, making computations more tractable \cite{oster22tensorHJB,dolgov2023datatensor}.

Amidst this landscape of approximation strategies, the State-Dependent Riccati Equation (SDRE) method has risen to prominence as a particularly effective and computationally appealing approach.  SDRE offers a pragmatic alternative by constructing local approximations of the optimal value function grounded in a Riccati equation framework.  The core idea behind SDRE is to extend the well-established Linear Quadratic Regulator (LQR) design methodology, renowned for its computational simplicity and guaranteed stability properties for linear systems, to the realm of nonlinear systems.  This extension is achieved by cleverly representing the nonlinear system dynamics in a "state-dependent linearized" form, also termed a "quasi-linear" or "state-dependent coefficient (SDC)" representation \cite{cliff1997sdre_survey,cimen2008nonlinearcontrol}.  This SDC transformation involves rewriting the nonlinear system equations in a linear-like structure where the system matrices are no longer constant but become functions of the system's current state.  Crucially, this reformulation allows one to solve a sequence of algebraic Riccati equations, parameterized by the state, at each time step or iteration to synthesize a state feedback control law.

The SDRE method elegantly bridges the gap between linear and nonlinear control design. By recasting the nonlinear problem into a sequence of linear-like problems, it inherits the computational efficiency and structured design process of LQR, while simultaneously adapting to the system's nonlinear behavior through the state-dependent system matrices. This adaptability allows SDRE to capture essential nonlinear dynamics, leading to improved performance compared to purely linear control designs applied to nonlinear systems.  A significant advantage of SDRE lies in its computational feasibility, especially in real-time control contexts, as solving algebraic Riccati equations is numerically well-conditioned and computationally efficient compared to solving the full HJB PDE \cite{feron2012realtimeSDRE}. While SDRE solutions are generally suboptimal compared to the true HJB solution, they often provide a favorable trade-off between performance and computational cost, making them attractive for practical applications.

The rigorous foundations of SDRE-based control were laid by Banks and his colleagues, who provided a comprehensive formulation and demonstrated its effectiveness for stabilizing nonlinear systems under certain regularity conditions \cite{banks2007sdre_theory}.  Their work established theoretical underpinnings and explored conditions for existence and stability of SDRE controllers. Cloutier and collaborators further propelled the method's practical adoption by developing structured numerical solvers and showcasing real-world applications, particularly in aerospace and missile guidance \cite{cloutier1996sdre_applications,mracek1995fullauthoritySDRE}.  These contributions highlighted the practical applicability and computational advantages of SDRE.

More recent research endeavors have focused on enhancing the SDRE methodology along several key directions. Optimization-based SDRE formulations have been explored to systematically minimize approximation errors inherent in the SDC representation and improve the resulting control performance. For example, Jones et al. investigated optimizing SDC parameterizations to minimize a quadratic error functional related to the SDRE approximation \cite{jones2020optimizingSDRE}, while Dolgov et al. explored optimizing SDC decompositions within a tensor framework for improved computational efficiency and accuracy \cite{dolgov2022optimizingSDC}. Despite these advancements, the SDRE approach is not without its limitations. A critical aspect affecting the solution quality is the inherent non-uniqueness of the SDC representation.  For a given nonlinear system, multiple valid SDC decompositions often exist, and the choice of decomposition can significantly impact the accuracy, stability, and robustness of the resulting SDRE controller \cite{cimen2001uniquenessSDC,cliff1997sdre_survey}.  Existing methods for selecting an "appropriate" or "optimal" SDC decomposition often rely on heuristics, engineering intuition, or problem-specific insights, lacking a universally applicable systematic approach.  Furthermore, while SDRE offers computational advantages over direct HJB solution, the iterative solution of state-dependent Riccati equations still poses numerical challenges, particularly for high-dimensional systems or when real-time performance is paramount.  Efficient numerical solvers are thus crucial for practical SDRE implementation.  Prominent approaches include iterative methods like the Newton-Kleinman algorithm \cite{kleinman1968newtonkleinman} for solving algebraic Riccati equations, known for its quadratic convergence, and offline-online strategies, such as those proposed by Palitta et al., which pre-compute parts of the solution offline to accelerate real-time computation \cite{palitta2018offlineonlineSDRE}. Each numerical method presents its own set of trade-offs in terms of computational efficiency, memory requirements, and solution accuracy, necessitating careful consideration based on the specific application requirements.

However, even with the computational advantages of SDRE, challenges remain, particularly for high-dimensional systems or infinite-dimensional systems described by PDEs, necessitating efficient numerical solvers and potentially model order reduction techniques to maintain computational tractability \cite{Heiland_Werner_LPV_SDRE_NavierStokes, Alla_Kalise_Simoncini_SDRE_PDE}.


This paper aims to delve deeper into the theoretical underpinnings and practical aspects of the SDRE method. We provide a comprehensive analysis of its theoretical properties, focusing on deriving quantitative error bounds for the SDRE approximation to the optimal HJB solution. We also investigate systematic approaches for obtaining improved SDC representations that minimize the inherent approximation error and enhance control performance.  Furthermore, we explore and compare the computational characteristics of different numerical techniques for efficiently solving the sequence of Riccati equations arising within the SDRE framework, with a particular focus on iterative methods.  Our key contributions are summarized as follows:

\begin{itemize}
\item Error Bound Derivation: We derive novel error bounds for the SDRE approximation, providing a quantitative measure of the deviation between the suboptimal SDRE solution and the true optimal solution obtained from the HJB equation. These bounds offer insights into the accuracy and limitations of the SDRE approach and provide guidance for assessing its applicability in different scenarios.
\item Optimal Semilinear Representations: We investigate the problem of selecting optimal semilinear (SDC) representations. We propose and analyze systematic methodologies for minimizing the SDRE residual error. This includes formulating optimization problems to find SDC decompositions that minimize a relevant error functional, leading to improved approximation accuracy and potentially enhanced closed-loop performance.
\item Comparative Study of Numerical Solvers: We conduct a comparative study of prominent numerical methods for solving the sequence of algebraic Riccati equations within the SDRE framework.  This includes a detailed evaluation of iterative methods like the Newton-Kleinman algorithm and offline-online strategies, highlighting their respective advantages and disadvantages in terms of computational efficiency, convergence properties, and suitability for different application contexts. We illustrate our findings through numerical experiments, including a challenging application involving the control of a nonlinear reaction-diffusion PDE, showcasing the practical implications of our analysis.
\end{itemize}

The remainder of this paper is structured as follows: Section 2 formally states the infinite-horizon nonlinear optimal control problem and establishes its fundamental connection to the Hamilton-Jacobi-Bellman equation. Section 3 rigorously introduces the State-Dependent Riccati Equation approach, detailing the SDC transformation and the derivation of the SDRE control law, and explores its key theoretical properties. Section 4 is dedicated to numerical techniques for solving the SDRE, including a detailed comparative analysis of two distinct computational methods, exemplified through the control of a nonlinear reaction-diffusion PDE. Finally, Section 5 concludes the paper with a summary of our key findings and a discussion on potential avenues for future research and extensions of the SDRE methodology.

Optimal control of nonlinear dynamical systems is a fundamental challenge in engineering, economics, and applied mathematics. The Hamilton-Jacobi-Bellman (HJB) equation provides a rigorous framework for computing optimal feedback control laws, stemming from Bellman’s groundbreaking work on dynamic programming \cite{bellman1957dynamic}. By dictating the evolution of the optimal value function through Bellman’s principle of optimality, the HJB equation yields the desired control policy via differentiation. Despite its theoretical elegance and completeness, directly solving the HJB equation is practically intractable. Its inherent nonlinearity and the exponential increase in computational complexity with the state-space dimension—the so-called “curse of dimensionality” \cite{bellman2015dynamic}—render exact solutions unattainable for most realistic systems, whether due to the analytical difficulty of a nonlinear partial differential equation (PDE) or the prohibitive cost of numerical discretization.

To overcome these challenges, a wide array of approximation techniques have been developed. Methods based on sparse grid discretizations \cite{GK16sparsegrids} alleviate the curse of dimensionality in numerical PDE solvers, while approaches employing artificial neural networks (ANNs) \cite{Darbon_Langlois_Meng_2020,Kunisch_Walter_2021,Zhou_2021} and regression-type methods using tensor formats \cite{oster22,dolgov2023data} offer alternative routes to approximate the value function and control policy. Neural network techniques, for instance, parameterize the value function and optimize network weights to minimize the residual of the HJB equation, making them powerful tools for high-dimensional problems. Tensor-based methods similarly aim to represent high-dimensional functions in compressed formats, facilitating efficient computation.

Amidst this landscape, the State-Dependent Riccati Equation (SDRE) method has emerged as a particularly effective and computationally feasible alternative. The core idea of SDRE is to extend the classical Linear Quadratic Regulator (LQR) design—which is known for its computational simplicity and robust stability for linear systems—to the nonlinear domain by recasting the system dynamics into a state-dependent linearized form. This “quasi-linear” or “state-dependent coefficient (SDC)” representation \cite{cliff1997sdre_survey,cimen2008nonlinearcontrol} rewrites the nonlinear equations so that the system matrices become functions of the current state. As a result, a sequence of algebraic Riccati equations can be solved, each parameterized by the state, to yield a local approximation of the optimal value function. Although SDRE solutions are generally suboptimal compared to the full HJB solution, they offer a favorable trade-off between control performance and computational efficiency—especially in real-time applications \cite{feron2012}.

The rigorous foundations of SDRE-based control were laid by Banks et al. \cite{banks2007}, who demonstrated its effectiveness in stabilizing nonlinear systems under mild regularity conditions. Cloutier and collaborators \cite{cloutier1996} further advanced the method by introducing structured numerical solvers suitable for real-time implementation, particularly in aerospace and missile guidance applications \cite{mracek1995fullauthoritySDRE}. More recent studies have focused on enhancing the SDRE framework through optimization-based formulations \cite{jones2020,dolgov2022optimizing} and systematic strategies for selecting the optimal SDC representation. The choice of semilinear decomposition is critical, as different representations can significantly impact the accuracy, stability, and robustness of the resulting SDRE controller. Moreover, the iterative solution of state-dependent Riccati equations—employing methods such as the Newton-Kleinman algorithm \cite{kleinman1968} and offline-online strategies \cite{palitta2018}—introduces additional numerical challenges, particularly for high-dimensional or infinite-dimensional systems described by PDEs.

In this paper, we delve deeper into the theoretical underpinnings and practical aspects of the SDRE method. Our main contributions are as follows:

Error Bound Derivation: We derive novel error bounds for the SDRE approximation, quantitatively measuring the deviation between the suboptimal SDRE solution and the optimal solution prescribed by the HJB equation. These bounds provide critical insights into the accuracy and limitations of the SDRE approach.

Optimal Semilinear Representations: We investigate systematic methodologies for selecting improved semilinear (SDC) representations that minimize the residual error inherent in the SDRE approximation. This includes formulating optimization problems aimed at identifying SDC decompositions that enhance approximation accuracy and closed-loop performance.

Comparative Study of Numerical Solvers: We conduct a detailed comparative analysis of numerical methods for solving the sequence of Riccati equations inherent to the SDRE framework. Our study evaluates iterative approaches, such as the Newton-Kleinman algorithm, alongside offline-online strategies, highlighting their respective trade-offs in computational efficiency, convergence properties, and applicability to various control scenarios.

The remainder of this paper is structured as follows: Section 2 formally states the infinite-horizon nonlinear optimal control problem and establishes its connection to the HJB equation. Section 3 rigorously introduces the SDRE approach, detailing the SDC transformation and the derivation of the SDRE control law, along with an exploration of its theoretical properties. Section 4 is dedicated to numerical techniques for solving the SDRE, including a comparative analysis of two distinct computational methods demonstrated through the control of a nonlinear reaction-diffusion PDE. Finally, Section 5 concludes the paper with a summary of our key findings and a discussion of potential avenues for future research in SDRE-based control.

\end{comment}




\section{The Infinite Horizon Optimal Control Problem}\label{sec:iocp}
In this section, we establish the formulation of the deterministic infinite-horizon optimal control problem, for which we aim to derive a feedback control law. We begin by presenting the optimal feedback synthesis via the Hamilton-Jacobi-Bellman framework.  

We consider a control-affine dynamical system governed by  
\begin{equation}\label{eq}
\left\{ \begin{array}{l}
\dot{y}(s) = f(y(s)) + B(y(s)) u(s), \quad s \in (0, +\infty),\\
y(0) = x \in \mathbb{R}^d.
\end{array} \right.
\end{equation}  
Here, $y: [0, +\infty) \to \mathbb{R}^d$ represents the system state, while $u: [0, +\infty) \to \mathbb{R}^m$ denotes the control input. The set of admissible control functions is given by \second{$\mathcal{U} = L^\infty ([0, +\infty); \mathbb{R}^m)$}. The system dynamics are described by the continuously differentiable functions $f: \mathbb{R}^d \to \mathbb{R}^d$ and $B: \mathbb{R}^d \to \mathbb{R}^{d \times m}$, satisfying $f({0})  = {0}$. When emphasizing the dependence of the control input on the initial state $x \in \mathbb{R}^d$, we explicitly denote it as $u(t;x)$.  

The objective is to minimize the following infinite-horizon cost functional:  
\begin{equation}\label{cost}
 J(u(\cdot \, ;x)) := \int_0^{+\infty} y(s)^{\top} Q y(s) + u^{\top}(s) R u(s)\,ds\,,
\end{equation}  
where $Q \in \mathbb{R}^{d \times d}$ is a symmetric positive semidefinite matrix, and $R \in \mathbb{R}^{m \times m}$ is a symmetric positive definite matrix. The goal is to determine an optimal control policy in feedback form, meaning that the control law depends solely on the current state of the system.  

To this end, we define the value function corresponding to a given initial condition $x \in \mathbb{R}^d$ as  
\begin{equation}
V(x) := \inf\limits_{u\in\mathcal{U}} J(u(\cdot;x))\,,
\label{VF}
\end{equation}  
which, according to standard dynamic programming principles, satisfies the HJB PDE for every $x \in \mathbb{R}^d$:  
\begin{equation}\label{HJB}
\min\limits_{u\in \second{\mathbb{R}^m} }\left\{(f(x)+B(x)u)^{\top} \nabla V(x) + x^{\top} Q x+ u^{\top} R u \right\}=0.
\end{equation}  

The HJB equation \eqref{HJB} is a first-order, fully nonlinear PDE defined over $\mathbb{R}^d$, making it computationally challenging, particularly when $d$ is large.
\second{We remark that the finiteness of the value function is not guaranteed a priori in nonlinear infinite-horizon problems. However, as shown in the subsequent section, under suitable assumptions on the data, the SDRE feedback yields a locally exponentially stable closed-loop system, from which it follows that the associated infinite-horizon cost is finite for all initial conditions in the corresponding region of attraction.
}

%In the unconstrained setting where $U = \mathbb{R}^m$, 
The optimal control minimizing the left-hand side of \eqref{HJB} can be derived explicitly as  
\begin{equation}\label{optc}
u^*(x) = -\frac{1}{2} R^{-1} B(x)^{\top} \nabla V(x)\,,
\end{equation}  
leading to the following unconstrained form of the HJB equation:  
\begin{equation}\label{hjbu}
\nabla V(x)^{\top} f(x) - \frac{1}{4} \nabla V(x)^{\top} B(x) R^{-1} B(x)^{\top} \nabla V(x) + x^{\top} Q x = 0.
\end{equation}  

The solution of the partial differential equation \eqref{hjbu} presents significant computational challenges due to the well-known \emph{curse of dimensionality}. Specifically, on a structured grid with $N$ grid points per dimension, the computational complexity scales as $O(N^d)$, rendering the problem intractable for high-dimensional systems. Rather than directly computing the optimal value function and the associated optimal trajectory, one may relax these requirements by seeking a \emph{suboptimal} yet stabilizing feedback control strategy. This consideration has contributed to the widespread adoption of the SDRE approach. In the following section, we introduce its formulation and examine its relationship with the HJB equation.


\subsection{State-Dependent Riccati Equation}
\label{sec:SDRE}
\both{Under the regularity assumptions introduced in the previous section, Proposition~1 in~\cite{cimen2010systematic} guarantees the existence of a state-dependent matrix function $A(x): \mathbb{R}^{d} \rightarrow \mathbb{R}^{d \times d}$ such that $f(x)=A(x)x$. Accordingly, the dynamical system \eqref{eq} can be written in the semilinear form}

%Consider a dynamical system represented in semilinear form as

\begin{equation}
	\dot{y}(t)  = A(y(t)) y(t) +B(y(t)) u(t).
	\label{semilinear}
\end{equation}
\both{
In dimension greater than one, this representation is non-unique, a feature that will be exploited in the following discussion. Detailed results concerning uniqueness can be found in \cite{cloutier1996nonlinear}.
}
In the special case where the matrix functions are constant, i.e., $A(y(t)) = A \in \mathbb{R}^{d\times d}$ and  $B(y(t)) = B \in \mathbb{R}^{d\times m}$, the problem reduces to the well-known Linear Quadratic Regulator.
\begin{comment}
In this framework, we present the concepts of stabilizability and detectability, which will be fundamental to the forthcoming analysis.
\begin{definition}
    The pair $(A,B)$ is stabilizable if there exists
a feedback matrix $K \in \mathbb{R}^{m \times d}$ such that $A -B K$
is stable, i.e., if all its eigenvalues lie on the
open left half complex plane. 
\end{definition}


\begin{definition}
    The pair $(A,C)$ is detectable if the pair $(A^\top, C^\top)$ is stabilizable.

\end{definition}

\end{comment}
\second{

To extend the classical LQR framework to the state-dependent setting, we introduce stabilizability and detectability for non-constant matrix functions defined on a set $\Omega \subseteq \mathbb{R}^d$.

\begin{definition}
    Let $\Omega \subseteq \mathbb{R}^d$ be a nonempty set.  
    The pair $(A(\cdot), B(\cdot))$ is said to be stabilizable on $\Omega$ if for every $x \in \Omega$ there exists a feedback matrix $K(x) \in \mathbb{R}^{m \times d}$  
    such that the matrix $A(x) - B(x) K(x)$ is stable, that is, all its eigenvalues lie strictly in the open left half of the complex plane.
\end{definition}

\begin{definition}
    Let $\Omega \subseteq \mathbb{R}^d$ be a nonempty set.  
    The pair $(A(\cdot), C)$ is said to be detectable on $\Omega$ if the pair $(A^\top(\cdot), C^\top)$ is stabilizable on $\Omega$ 
    in the sense of the previous definition.
\end{definition}

When $A(\cdot)$ and $B(\cdot)$ are constant, these definitions reduce to the classical LQR notions of stabilizability and detectability.} In this setting, when the pair $(A,B)$ is stabilizable and the pair $(A,Q^{1/2})$ is detectable, the optimal feedback control for the LQR is given by:
\begin{equation*}
	u(y) = -R^{-1} B^{\top} P y,
\end{equation*}
where $P\in\mathbb{R}^{d\times d}$ is the unique positive definite solution of the Continuous Algebraic Riccati Equation (CARE)
\begin{align*}
	A^\top P + P A
	-P BR^{-1}B^\top P +Q = 0\ .
	%\label{are}
\end{align*}
The SDRE method generalizes this approach by allowing the system matrices to depend on the state, leading to the control law
\begin{equation}
	u_S(y) = -R^{-1} B^{\top}(y) P(y)y\,,
	\label{control_sdre}
\end{equation}
where $P(y)$ now satisfies a State-Dependent Riccati Equation:
\begin{align}
	A^\top(y) P(y) +  P(y) A(y)-P(y)B(y)R^{-1}B^\top(y)P(y)+Q & = 0,
	\label{sdre}
\end{align}
with the matrices $A(y)$ and $B(y)$ evaluated at the state $y$. A practical implementation of SDRE-based control for nonlinear stabilization employs a receding horizon strategy. In Algorithm \ref{alg:sdre} the method is sketched. Specifically, at each discrete state \(y_S(t_k)\), the matrices in equation \eqref{sdre} are fixed, and the corresponding Riccati equation is solved for 
\( P(y_S(t_k))\). The control input  \(u_S(y_S(t_k))\) from equation \eqref{control_sdre} is then applied over a short horizon to evolve the system dynamics. This process is repeated at the next state \(y_S(t_{k+1})\) iteratively. 

\begin{algorithm}[htbp]
\caption{The reduced horizon-SDRE method}
\label{alg:sdre}{
\begin{algorithmic}[1]
\Statex{\textbf{Input:} $A(\cdot), B(\cdot), Q,R$, initial conditions $ y_0$, time discretization $\{t_i\}_{i=0}^{N_T}$ }
\Statex{\textbf{Output:} Controlled trajectory $\{y_S(t_i)\}_{i=0}^{N_T}$}

\State{$y_S(t_0):= y_0$}
\For{$n = 0, \dots, N_T-1$}

\State{Compute $P(y_S(t_{n}))$ from \eqref{sdre}}
\State{Build the control $u_S(y_S(t_{n}))$ from \eqref{control_sdre}}
\State{Integrate the system \eqref{semilinear} to obtain $y_S(t_{n+1})$}       
\EndFor 
\end{algorithmic}}
\end{algorithm}

Under suitable stability assumptions, it can be shown that the closed-loop system governed by the feedback law \eqref{control_sdre} exhibits local asymptotic stability, as formalized in the following theorem from \cite{banks2007}.

\begin{theorem}
    Assume that the system \eqref{eq} is such that \( f(x) \) and \( \frac{\partial f(x)}{\partial x_j} \) (for \( j = 1, \dots, n \)) are continuous in \( x \) for all \( \|x\| \leq \hat{r} \), where \( \hat{r} > 0 \).
Assume further that \( A(x) \) and \( B(x) \) are continuous and that the pair $(A(x),B(x))$ is stabilizable and the pair $(A(x),Q^{1/2})$ is detectable in some nonempty neighborhood of the origin \( \mathcal{O} \subseteq B_{\hat{r}}(0) \). Then, the system with the control given by \eqref{control_sdre} is locally asymptotically stable.
\label{thm:stable}
\end{theorem}

\second{Indeed, the authors obtain a stronger result. The system controlled by the feedback law \eqref{control_sdre} is locally exponentially stable, \emph{i.e.}, there exist constants $r>0$, $C>0$, and $\lambda>0$ such that
$\|y_S(t)\| \le C \|y_S(0)\| e^{-\lambda t}$, for all $\|y_S(0)\|\le r,\; t \in [0,+\infty)$.
This exponential decay implies that the quadratic running cost is integrable over $[0,+\infty)$, and therefore the value function is finite for all initial states within the region of attraction of the stabilizing feedback, as anticipated in Section~\ref{sec:iocp}.} 

Beyond stabilization, an alternative perspective considers the deviation of the SDRE approach from the optimality conditions dictated by HJB equation. To analyze this discrepancy, we introduce the SDRE-based approximation of the value function:
\begin{equation}
    V_{S}(x) = x^\top P(x) x,
    \label{v_sdre}
\end{equation}  
where \( P(x) \) satisfies the SDRE \eqref{sdre}. The feedback control \eqref{control_sdre} is derived from the optimal control formula \eqref{optc}, under the assumption that the term involving the gradient of \( P(x) \) is negligible.
However, accounting for this term leads to the modified feedback law:

\begin{equation}
\tilde{u}_{S}(x) = -\frac{1}{2} R^{-1} B(x)^T \nabla V_S (x),
\label{utilde_sdre}
\end{equation}
where $\nabla V_S (x) = 2P(x) x + \varphi(x)$,
with the correction term $\varphi$ defined as
$$
[\varphi(x)]_i =  x^\top P_{x_i}(x) x, \quad i =1, \dots, d,
$$
with \( P_{x_i}(x) = \frac{\partial P(x)}{\partial x_i} \). By employing similar analytical techniques as those utilized in the proof of Theorem \ref{thm:stable} in \cite{banks2007}, one can establish that the system dynamics under the control $\tilde{u}_{S}(x)$ exhibit local asymptotic stability, as formally stated in the following result.


\begin{theorem}

% Assume that the system \eqref{eq} is such that \( f(x) \) and \( \frac{\partial f(x)}{\partial x_j} \) (for \( j = 1, \dots, n \)) are continuous in \( x \) for all \( \|x\| \leq \hat{r} \), where \( \hat{r} > 0 \). Assume further that \( A(x) \) and \( B(x) \) are continuous and that the pair $(A(x),B(x))$ is stabilizable and the pair $(A(x),Q^{1/2})$ is detectable in some nonempty neighborhood of the origin \( \mathcal{O} \subseteq B_{\hat{r}}(0) \). Then, the system with the control given by \eqref{utilde_sdre} is locally asymptotically stable.
Under the assumptions of Theorem \ref{thm:stable}, the system with the control given by \eqref{utilde_sdre} is locally \second{exponentially} stable.
\label{thm:stable2}
\end{theorem}

\begin{proof}
    Under the feedback control
$\tilde{u}_{S}(x)$, the closed-loop dynamics becomes
$$
    \dot{y} = A(y)y - B(y) R^{-1} B(y)^T P(y)y - \frac{1}{2} B(y) R^{-1} B(y)^T \varphi(y).
    \label{closed_loop}
$$
Equivalently, this can be expressed as
$$
 \dot{y} = (A_0-B_0 R^{-1}B_0^\top P_0)y + h(y),
$$
where $A_0=A(0)$, $B_0 = B(0)$, $P_0 = P(0)$ and 
\second{
\begin{equation*}
\begin{aligned}
h(y)
&= \bigl(A(y) - A_0\bigr)\,y
   - \Bigl(B(y) R^{-1} B(y)^\top P(y) - B_0 R^{-1} B_0^\top P_0\Bigr)\,y \\
&\quad - \frac{1}{2}\, B(y) R^{-1} B(y)^\top \varphi(y).
\end{aligned}
\end{equation*}
}


%$h(y)$ represents a perturbation term.
It follows directly that
$$
    \lim_{y \to 0} \frac{\|h(y)\|}{\|y\|} = 0,
$$
since all the state-dependent matrices in $h(y)$ are continuous and the additional term involving $\varphi$ constitutes a perturbation of higher order in $y$. As a result, the stability analysis proceeds within the same framework as employed in the proof of Theorem \ref{thm:stable} from \cite{banks2007}.
 
\end{proof}



Now, given the SDRE ansatz \eqref{v_sdre}, it is possible to compute its residual with respect to the HJB equation \eqref{hjbu} just inserting $\nabla V_S(x)$ into \eqref{hjbu}, obtaining 
$$
    x^\top \big( P(x) A(x) + A^\top(x) P(x) - P(x) B(x) R^{-1} B^\top(x) P(x) + Q \big) x +E(x) = 0,
    $$
    with
    \begin{equation}
        E(x) = \varphi(x)^\top   \left([A(x)- B(x)R^{-1}B(x)^\top P(x)] x - \frac{1}{4} B(x) R^{-1} B(x)^\top \varphi(x) \right).
        \label{residual}
    \end{equation}
The first term in the expression corresponds to the SDRE \eqref{sdre}. As a result, it is either exactly zero in the theoretical formulation or negligibly small when considering numerical approximations, depending on the accuracy of the numerical method employed. The function $E(x)$ quantifies the effect of the sub-optimality associated with the SDRE approach. While the feedback control $u_S$ provides a stabilizing suboptimal control, the function $\tilde{u}_{S}$ accounts for the correction term $E(x)$, making it a refined alternative when optimality is a concern. Indeed, minimizing $E(x)$  aligns the feedback law $\tilde{u}_{S}$ with the optimal control solution derived from the HJB equation. 
%The function $\tilde{u}_{S}$ accounts for this correction, making it a refined alternative when optimality is a concern.
We note that fixing the semilinear form $A(x)$, the residual $E(x)$ can be explicitly determined by computing the derivatives \( \{P_{x_i}(x)\}_i \). Whenever we want to stress the dependence on the semilinear form, we write $P(x;A(x))$ and $E(x;A(x))$.

To obtain the derivatives \( \{P_{x_i}(x)\}_i \), we differentiate the SDRE \eqref{sdre} with respect to \(x_i\), yielding the following Lyapunov equation:

\begin{equation}
P_{x_i}(x)A_{cl}(x) + A_{cl}^\top(x)P_{x_i}(x) 
+\Lambda_i(x) = 0,\label{Lyap}
\end{equation}
where 
\begin{align*}
\Lambda_i(x) =\, &P(x)A_{x_i}(x) + A^\top_{x_i}(x)P(x) \nonumber\\[1mm]
& - P(x) (B_{x_i}(x)R^{-1}B^\top(x) + B(x)R^{-1}B^\top_{x_i}(x))P(x)
\end{align*}
and
\[
A_{cl}(x) = A(x) - B(x)R^{-1}B^\top(x)P(x),
\]
which is referred to as the \emph{closed-loop} system matrix.


Equation \eqref{Lyap} defines a Lyapunov equation for \(P_{x_i}(x)\). Solving this equation for each $i \in \{1, \dots, d\}$ allows for the computation of the residual term \eqref{residual}, thereby quantifying the deviation introduced by the omitted gradient terms.
In particular, if the residual vanishes for all $x$, then the SDRE formulation is equivalent to the HJB equation. Conversely, if the residual does not vanish, one may incorporate $E(x)$ into the running cost to show that the function $V_S(x)$ satisfies a modified HJB equation. Consequently, $V_S(x)$ adheres to a dynamic programming principle, as formalized in the following result.


    

\begin{comment}
    

Even if the previous results ensures the local asymptotic stability, one may be interested in how this feedback law differs from the one obtained via the HJB formulation. For this reason, we introduce the following SDRE approximation of the value function
\begin{equation}
    V_{\text{SDRE}}(x) = x^\top P(x) x,
    \label{v_sdre}
\end{equation}  
where \( P(x) \) satisfies the SDRE \eqref{sdre}. The corresponding feedback control \( u_{\text{SDRE}} \) is derived from the optimal control formula \eqref{optc}, under the assumption that the term involving the gradient of \( P(x) \) is negligible.  

Explicitly, the gradient of the approximation \eqref{v_sdre} is given by  
\begin{equation}
    \nabla V_{\text{SDRE}}(x) = 2P(x)x + \varphi(x),
    \label{nabla_V_sdre}
\end{equation}  
where
$$
[\varphi(x)]_i =  x^\top P_{x_i}(x) x, \quad i =1, \dots, d,
$$
with \( P_{x_i}(x) = \frac{\partial P(x)}{\partial x_i} \).% The omission of the term involving the gradient of \( P(x) \) introduces an inherent suboptimality in the SDRE-based control strategy.


It is possible to compute the misfit between the HJB equation \eqref{hjbu} and the SDRE ansatz \eqref{v_sdre} just inserting \eqref{nabla_V_sdre} into \eqref{hjbu}, obtaining 
$$
    x^\top \big( P(x) A(x) + A^\top(x) P(x) - P(x) B(x) R^{-1} B^\top(x) P(x) + Q \big) x +E(x) = 0,
    $$
    with
    \begin{equation}
        E(x) = \varphi(x) \cdot \left(A(x)x - \frac{1}{4} B(x) R^{-1} B(x)^\top \varphi(x) \right).
        \label{residual}
    \end{equation}


The first term in the expression corresponds to the SDRE \eqref{sdre}. As a result, it is either exactly zero in the theoretical formulation or negligibly small when considering numerical approximations, depending on the accuracy of the numerical method employed.

The function $E(x)$ quantifies the effect of the omitted gradient terms, which account for the sub-optimality associated with the SDRE approach. 
%Specifically, if $E(x)\le 0$, then $V_S(x)$ serves as a \emph{supersolution} for the HJB equation, while if $E(x)\ge 0$ it constitutes a \emph{subsolution}. For a detailed discussion on the concepts of sub- and supersolutions, we refer to \cite{BCD97}. 
\end{comment}



%If the residual is null for every $x$, we have the equivalence between the SDRE formulation and the HJB equation. In the case it is not null, we note that incorporating the $E(x)$ into the running cost, we can estalish that the functin $V_S(x)$ satisfies a \emph{modified} HJB equation and, hence, satisfies a dynamic programmin principle, as stated in the following result. 



\begin{proposition}
Define the augmented running cost by
\begin{equation}\label{eq:augmented_cost}
\tilde{\ell}(x,u) = x^\top Q x + u^\top R u + E(x),
\end{equation}
where $E(x)$ is given by \eqref{residual}.
Then, the function $V_S(x)$
satisfies:
\begin{itemize}
    \item The modified HJB equation:
    \begin{equation}
    \min_{u\in \mathbb{R}^m} \left\{ \nabla V_S(x)^\top\Bigl(A(x)x+B(x)u\Bigr) + \tilde{\ell}(x,u) \right\} = 0.
    \label{hjb_sdre}
        \end{equation}
    \item The dynamic programming principle (DPP):
       \begin{equation}
    V_S(x) = \inf_{u \in \mathcal{U}} \left\{ \int_0^T \tilde{\ell}\bigl(x(s),u(s)\bigr)\,ds + V_S\bigl(x(T)\bigr) \right\},
        \label{dpp_sdre}
        \end{equation}
    for every \(T\ge 0\).
\end{itemize}
\end{proposition}

\begin{proof}
As $E(x)$ is defined as the residual of $V_S(x)$ with respect to the HJB equation \eqref{HJB}, it satisfies a modified HJB equation in which the running cost is given by the augmented formulation \eqref{eq:augmented_cost}. The DPP formulation in \eqref{dpp_sdre} follows directly from the classical theory of HJB equations; see, for instance, \cite{BCD97} for further details.
 
\end{proof}

\begin{example}[Linear Quadratic Regulator]
A fundamental and straightforward example is the LQR. In this case, the system matrices $A(x)=A$ and $B(x)=B$ are constant, implying that the solution to the SDRE is also constant. Specifically, under these conditions, the HJB equation and the SDRE approach (which, in this case, reduces to the standard LQR formulation) coincide, yielding the exact optimal value function and the corresponding optimal feedback control law.
    
\end{example}

\begin{example}[Van Der Pol oscillator]
\label{ex:vdp}
We consider the dynamics

\[
\dot{x}=A(x)x+B(x)u,
\]
with
\[
A(x)=\begin{pmatrix}0 & 1\\[1mm]-1 & -0.5\,(1-x_1^2)\end{pmatrix},\quad B(x)=\begin{pmatrix}0\\ x_1\end{pmatrix}.
\]
and cost functional \eqref{cost} with
\[
Q=\begin{pmatrix}0 & 0\\[1mm]0 & 1\end{pmatrix},\quad R=1.
\]


A direct computation shows that the constant choice
\[
P(x)=I=\begin{pmatrix}1&0\\[1mm]0&1\end{pmatrix}
\]
satisfies the SDRE \eqref{sdre}. 

Since \(P(x)\) is constant, its derivatives vanish and the residual (which involves the derivatives $\{P_{x_i}(x)\}_i$) is identically zero. In this scenario, despite the state-dependent nature of the system matrices, full optimality is achieved as a consequence of the SDRE solution remaining constant.


\end{example}

















\section{Error Bounds based on the Residual}


We have established that the residual of the SDRE solution with respect to the HJB equation can be quantified using equation \eqref{residual}, which relies on solving the Lyapunov equations for the derivatives $\{P_{x_i}\}_i$. In this section, our objective is to derive an error bound between the SDRE approximation \eqref{v_sdre} and the optimal value function \eqref{VF}, based on the residual $E(x)$. For a discounted optimal control problem, where the cost functional is defined as
\begin{equation*}\label{cost2}
 J(u(\cdot \,;x)) := \int_0^{+\infty} e^{-\rho t}\left(y(s)^{\top} Q y(s) + u^{\top}(s) R u(s)\right) \,ds\,,
\end{equation*} 

with $\rho>0$, it follows that the dynamic programming operator exhibits $L^\infty$-contractive (see, $e.g.$, \cite{dolcetta1983discrete}). This property enables the derivation of an error estimate of the form
$$
\Vert \tilde{V} - V \Vert_{\infty} \le \frac{\Vert E_{\tilde{V}} \Vert_{\infty}}{\rho},
$$
where $\tilde{V}$ represents an approximation of the value function, and $E_{\tilde{V}}$ denotes its residual with respect to the HJB equation (see \cite{grune1997adaptive,albert2002posteriori} for related approaches). However, this methodology is not applicable in our setting, as we consider an undiscounted framework where $\rho = 0$. Instead, we employ the dynamic programming principle \eqref{dpp_sdre} for the perturbed value function $V_S(x)$ in conjunction with the local asymptotic stability of the SDRE-controlled trajectory. A similar approach has been explored in the extended abstract \cite{grune2018hamiltonian}, where the authors establish a qualitative error bound based on the residual analysis.

\begin{comment}
Furthermore, the SDRE approach relies on solving distinct Riccati equations at different state positions. As a result, the effective cost functional is obtained through time integration, defined as
\begin{equation}
    J_{SDRE}(x) := \int_{0}^{+\infty} y_S(s;x)^\top Q \, y_S(s;x) + u_S(s)^\top R \, u_S(s) \, ds,
\end{equation}
where $u_S$ is given by \eqref{control_sdre}, and $y_S(s;x)$ denotes the solution of the state dynamics \eqref{semilinear} with control $u_S$ starting from the initial condition $x$. The following result formalizes the relationship between the function $V_S(x)$ and the effective cost $J_{SDRE}(x)$ by incorporating the integral of the residual along the SDRE-controlled trajectory. 

\begin{proposition}
Under the assumptions of Theorem \ref{thm:stable}, the following equality holds: 
\begin{equation}
V_S(x) \le J_{SDRE}(x) +\int_0^{\infty} E(y_S(t)) dt,
\end{equation}
in some nonempty neighborhood of the origin.
\end{proposition}

%\begin{proof}
By applying the DPP \eqref{dpp_sdre}, we obtain the inequality
$$
V_S(x) \le \int_0^T \tilde{l}(y_S(s),u_S(s))+ V_S(y_S(T)).
$$
Taking the limit as $T \rightarrow + \infty$, and noting that Theorem \ref{thm:stable} ensures the local asymptotic stability of the SDRE trajectory, we conclude that $\lim_{T \rightarrow + \infty} V(y_S(T)) =0$ by continuity. Finally, invoking the definition of the augmented cost in \eqref{eq:augmented_cost}, the proof is complete.

    The residual $E(x)$ can be formulated as
\[
E(x) = L(x,u_S) + \nabla V_S(x)^\top f(x,u_S),
\]
where
\[
L(x,u) = x^\top Qx + u^\top Ru.
\]
 The \emph{effective} SDRE cost is then defined as
\[
V^e_{SDRE}(x) := \int_0^{+\infty} L\bigl(y_S(t;x), u_S(t)\bigr)dt.
\]


Differentiating $V_S\bigl(y_S(t;x)\bigr)$ along the trajectory gives
\[
\frac{d}{dt}V_S\bigl(y_S(t;x)\bigr)
= \nabla V_S\bigl(y_S(t;x)\bigr)^\top f\bigl(y_S(t;x), u_S(t)\bigr).
\]
Thus,
\[
\frac{d}{dt}V_S\bigl(y_S(t;x)\bigr) + L\bigl(y_S(t;x), u_S(t)\bigr)
= E\bigl(y_S(t;x)\bigr).
\]
Integrating from $t=0$ to $T$,
\[
V_S\bigl(y_S(T;x)\bigr) - V_S(x)
+ \int_0^\top L\bigl(y_S(t;x), u_S(t)\bigr)dt
= \int_0^\top E\bigl(y_S(t;x)\bigr)dt.
\]
Assuming stability so that $\lim_{T\to\infty}V_S\bigl(y_S(T;x)\bigr)=0$, letting $T\to\infty$ gives
\[
V^e_{SDRE}(x) = V_S(x) + \int_0^{+\infty} E\bigl(y_S(t;x)\bigr)dt.
\]
\end{comment}
%\end{proof}
%\end{comment}
% The bound proposed in the current work 
%The aforementioned proof relies on the application of the DPP \eqref{dpp_sdre} for the perturbed value function $V_S(x)$ and the local asymptotic stability of the SDRE-controlled trajectory. The same principles hold for the optimal value function $V(x)$, enabling us to establish the following error bound, which is derived from the integration of the residual along both the optimal trajectory, computed using the optimal control, and the SDRE-controlled trajectory. Similar ideas have been used in \cite{grune2018hamiltonian} in the conte

\begin{proposition}
Under the assumptions of Theorem \ref{thm:stable}, the following error bound holds in some nonempty neighborhood of the origin
    \begin{equation}
|V_S(x) - V(x)| \leq \max\Bigl\{\int_0^{\infty} |E(y^*(t))| dt, \int_0^{\infty} |E(\tilde{y}_{S}(t))| dt\Bigr\}, 
\label{eq:bound}
\end{equation}
where $y^*(t)$ denotes the state trajectory corresponding to the optimal control \eqref{optc}, and $\tilde{y}_{S}(t)$  is the trajectory associated with the control \eqref{utilde_sdre}.
%\begin{equation}
%\tilde{u}=-\frac{1}{2} R^{-1} B(x)^\top \nabla V_S(x).
%\label{eq:utilde}
%\end{equation}

Furthermore, assume that both controlled trajectories are exponentially stable and remain in a neighborhood $\Omega$ of the origin. Then there exists a constant $C>0$ such that
\begin{equation}
|V_S(x)-V(x)|
\le
\frac{C}{\lambda}\,\|E\|_{*,\Omega}\,\|x\|,
\qquad x\in\Omega,
\label{eq:error_weighted}
\end{equation}
where
\[
\|E\|_{*,\Omega}:=\sup_{z\in\Omega\setminus\{0\}} \frac{|E(z)|}{\|z\|}.
\]



\end{proposition}

\begin{proof}


Applying the DPP for the optimal value function, we obtain that for any \(T\ge 0\)
\begin{equation}\label{eq:Vopt}
V(x)=\int_0^T \ell\bigl(y^*(s),u^*(s)\bigr)ds+V\bigl(y^*(T)\bigr),
\end{equation}
where \(y^*(t)\) is the state trajectory corresponding to the optimal control \(u^*(t)\).

Now, applying the same control \(u^*(\cdot)\) into the DPP for the SDRE approximation, we obtain
\begin{equation}\label{eq:VSDRE}
V_S(x)\le \int_0^T \Bigl[\ell\bigl(y^*(s),u^*(s)\bigr)+E\bigl(y^*(s)\bigr)\Bigr]ds+V_S\bigl(y^*(T)\bigr).
\end{equation}
By subtracting equation \eqref{eq:Vopt} from equation \eqref{eq:VSDRE} and utilizing the local asymptotic stability of the trajectory, the terminal terms vanish in the limit as $T \rightarrow + \infty$, leading to
\[
V_S(x)-V(x)\le \int_0^\infty E\bigl(y^*(s)\bigr)ds.
\]
%yields
%$$
%V_S(x)-V(x)\le  \int_0^T E\bigl(y^*(s)\bigr)ds+ \Bigl[V_S\bigl(y^*(T)\bigr)-V\bigl(y^*(T)\bigr)\Bigr].
%$$
%Since the trajectory is local asymptotically stable, the terminal terms vanish in the limit as $T \rightarrow + \infty$, leading to

A similar bound for $V(x)-V_S(x)$ can be derived by considering the control $\tilde{u}_{S}$ given in \eqref{utilde_sdre}, obtaining the bound \eqref{eq:bound}.

%Now, suppose further that the closed-loop systems under \( u^*(x) \) and $\tilde{u}(x)$ are exponentially stable. Then there exist a decay rate \( \lambda > 0 \) and a constant \( C > 0 \) such that along the trajectory,  
%\[
%E(y^*(t)) \leq C \|E\|_{\infty} e^{-\lambda t},
%\]
%so that  
%\[
% \int_0^\infty E(y^*(t)) \, dt  \leq \int_0^\infty C \|E\|_{\infty} e^{-\lambda t} \, dt = \frac{C}{\lambda} \|E\|_{\infty}.
%\]
%The same bound holds for the term involving $E(\tilde{y}_{S}(t))$, completing the proof.

Now, by definition of $\|E\|_{*,\Omega}$,
\[
|E(z)|\le \|E\|_{*,\Omega}\|z\|,
\qquad z\in\Omega.
\]
Hence, using exponential stability,
\[
|E(y^*(t;x))|
\le
\|E\|_{*,\Omega}\|y^*(t;x)\|
\le
C \|E\|_{*,\Omega} e^{-\lambda t}\|x\|,
\]
and similarly for $\tilde y_S(t;x)$. Therefore,
\[
\int_0^\infty |E(y^*(t;x))|\,dt
\le
C \|E\|_{*,\Omega}\|x\|
\int_0^\infty e^{-\lambda t}\,dt
=
\frac{C}{\lambda}\|E\|_{*,\Omega}\|x\|.
\]
The same estimate holds along $\tilde y_S(t;x)$, which proves
\eqref{eq:error_weighted}.

 
\end{proof}


The norm $\|E\|_{*,\Omega}$ in the error bound \eqref{eq:error_weighted} can be evaluated within a compact domain $\Omega$ provided that both the controlled trajectories $y^*(t)$ and $\tilde{y}_S(t)$ remain within $\Omega$.

The error bound given in \eqref{eq:bound} is purely theoretical, as it relies on knowledge of the optimal trajectory $y^*(t)$. However, if $V_S(x) \le V(x)$, the inequality is obtained considering the integral of the residual along the trajectory computed via $\tilde{u}_{S}$ and it is a practical indicator, since it is fully computable. This approach involves computing the controlled trajectory $\tilde{y}_{S}(t)$ and evaluating the residual based on the methodology outlined in Section \ref{sec:SDRE}.  In the following example, we show that, for the given scenario, this indicator, derived from the SDRE trajectory, guarantees the validity of the upper bound.






\begin{comment}
In many practical cases, one can estimate $E(x)$ as $\|E(x)\| \leq C \|x\|^3$, and if the closed-loop system decays exponentially as $\|x(t)\| \leq K e^{-\alpha t} \|x\|$, then
\begin{equation}
\int_0^{\infty} \|E(x(t))\| dt \leq \frac{C K^3}{3\alpha} \|x\|^3.
\end{equation}
This gives an explicit local error bound:
\begin{equation}
\|V_S(x) - V^*(x)\| \leq \frac{C K^3}{3\alpha} \|x\|^3.
\end{equation}


We assume:
\begin{itemize}
    \item The functions $P(x)$, $A(x)$, $B(x)$, and $R(x)$ are continuously differentiable in a neighborhood $\Omega$ of the origin.
    \item There exists a constant $M > 0$ such that
    \begin{equation}
        \|P_{x_i}(x)\| \leq M, \quad \forall x \in \Omega, \quad i = 1, \dots, n.
    \end{equation}
    \item The function 
    \begin{equation}
        F(x) := A(x)x - \frac{1}{4} B(x) R^{-1} B(x)^\top \varphi(x)
    \end{equation}
    is locally Lipschitz and satisfies the bound
    \begin{equation}
        \|F(x)\| \leq L\|x\|
    \end{equation}
    for some constant $L > 0$.
\end{itemize}
Then
$$
\| E(x) \|  \le \| \varphi(x) \| \| F(x) \| \le d M L \| x \|^3.
$$
\end{comment}

\begin{example}(Van Der Pol oscillator)

Let us consider again the setting of Example \ref{ex:vdp}. Let us consider a different choice for semilinear form
\first{
\[
A(x)=\begin{pmatrix} -x_2 & 1+x_1\\[1mm]-1 & -0.5\,(1-x_1^2)\end{pmatrix},\quad B(x)=\begin{pmatrix}0\\ x_1\end{pmatrix}.
\]
}
In this scenario, the semilinear form does not represent an optimal solution, and consequently, \(E(x)\) remains nonzero for certain values of $x\in \mathbb{R}^2$. Let us examine the initial condition \(x^0 = (-0.5, 0.5)\) and compute the trajectory governed by the controller $\tilde{u}_S$. Subsequently, we evaluate the error \(|V(x) - V_S(x)|\) along the SDRE-controlled trajectory, together with the error bound as specified in \eqref{eq:bound}. The results are presented in the left panel of Figure \ref{pics:vdp}. The integral of the residuals computed along the optimal trajectory \(y^*(t)\) and the trajectory \(\tilde{y}_{S}(t)\) exhibit similar behavior. Therefore, only the integral corresponding to the latter is presented in Figure \ref{pics:vdp}.
 It is observed that as the state approaches the origin, the error diminishes towards zero, while the error bound remains closely aligned with the peaks of the error, indicating its effectiveness as an error predictor. Finally, in the right panel of Figure \ref{pics:vdp}, the trajectory obtained through the SDRE controller is depicted, where it is evident that the system stabilizes towards the origin.


\begin{figure}[H]	
\begin{center}    
	\includegraphics[width=0.49\textwidth]{pics/error_bound.eps}	
 	\includegraphics[width=0.49\textwidth]{pics/traj_vdp.eps}	
\caption{Van Der Pol example.  \emph{Left}: The temporal evolution of both the error between the value function and its SDRE approximation and the corresponding error bound as given in \eqref{eq:bound} along the trajectory induced by the SDRE controller. \emph{Right}: controlled solution using the SDRE controller.}
\label{pics:vdp}
\end{center}


\end{figure}

\end{example}

\begin{comment}
    
\begin{example}[Cucker-Smale model]
Considering a system of $N_a$ interacting agents, we analyze the dynamics governed by the Cucker-Smale model, which is described by the following system:
\begin{equation*}
\begin{bmatrix} \dot{y} \\ \dot{v}  \end{bmatrix} = A(y)\begin{bmatrix} y \\ v  \end{bmatrix}+Bu= \begin{bmatrix} 0_{N_a} & I_{N_a} \\ 0_{N_a} & \mathcal{A}_{N_a}(y)  \end{bmatrix} \begin{bmatrix} y \\ v  \end{bmatrix} + \begin{bmatrix} O_{N_a} \\ I_{N_a}  \end{bmatrix} u,
\end{equation*}
where the interaction matrix is given by
\begin{equation*}
\left[ \mathcal{A}_{N_a}(y) \right]_{i,j}= \begin{cases} -\frac{1}{N_a} \sum\limits_{k \neq i} K(y_i,y_k), & \text{if } i=j, \\
\frac{1}{N_a} K(y_i,y_j), & \text{otherwise},
\end{cases}
\end{equation*}
with the interaction function
\begin{equation*}
K(y_i,y_j)=\frac{1}{1+\| y_i-y_j \|^2}.
\end{equation*}
Defining $y_v = [y; v] \in \mathbb{R}^n$, where $n=2N_a$, our objective is to regulate both positions and velocities to zero while minimizing the following cost functional:
\begin{equation*}
J(y(\cdot),v(\cdot),u(\cdot))=  \int_0^\infty  y(s)^\top Q_{N_a} y(s) + v(s)^\top  Q_{N_a} v(s)  +  u(s)^\top R u(s)\, ds,
\end{equation*}
where $Q_{N_a}=R=\frac{1}{N_a} I_{N_a} \in \mathbb{R}^{N_a \times N_a}$.

We consider an initial condition that is randomly distributed within the interval $[0,0.1].$
The controlled trajectory is computed using Algorithm \ref{alg:sdre}, where numerical integration is performed with the MATLAB function \texttt{ode45}, and the associated CAREs are solved using the function \texttt{icare}. The left panel of Figure \ref{pics:CS} illustrates the evolution of the residual $E(y_S(t))$ over time. It is observed that the residual remains consistently below $10^{-6}$ and decreases toward zero as the state approaches the origin. %In the subsequent section, we will demonstrate that the integral of the residual over the controlled trajectory serves as an upper bound for the optimal value function.
Given that the total cost is $0.0170$ and the integral of the residual over time is approximately $3 \cdot 10^{-7}$, the relative error can be estimated to be on the order of $O(10^{-5})$. The right panel of Figure \ref{pics:CS} depicts the controlled trajectory, illustrating the stabilizing effect of the SDRE controller on the system.




\begin{figure}[H]	
\begin{center}    
	\includegraphics[scale=0.38]{pics/CS/residual_CS.eps}	
    \includegraphics[scale=0.38]{pics/CS/traj_CS.eps}
\caption{Cucker Smale example. \emph{Left}: Evolution of the residual $E(x(t))$ along the time trajectory. \emph{Right}: trajectories controlled via the SDRE controller.}
\label{pics:CS}
\end{center}


\end{figure}



    
\end{example}
\end{comment}


\subsection{Optimal Semilinear Form}

In the previous sections, we introduced the SDRE methodology, which relies on the semilinear representation \eqref{semilinear} of the dynamical system. This representation plays a fundamental role, as different choices of semilinear forms can lead to varying outcomes in the control strategy.
A natural question that arises is whether there exists at least one semilinear decomposition for which the residual term $E(x)$ vanishes entirely. This is a critical consideration, as the presence of $E(x)$ contributes to the sub-optimality of the SDRE approach. The following result provides an affirmative answer to this question by establishing the existence of such a semilinear form, under the assumption that the residual exhibits a sign change for two distinct semilinear formulations.

\begin{proposition}
\first{Let us consider the assumptions of Theorem \ref{thm:stable} and a fixed $x\in\mathbb{R}^n$.}
 %   Let $f:\mathbb{R}^n\to\mathbb{R}^n$ be a smooth function and let $x\in\mathbb{R}^n$ be fixed. 
    Consider a fixed base semilinear form $A_0(x)$ satisfying $A_0(x)x=f(x)$, then the general semilinear form is parameterized as
\[
A(x)=A_0(x)+Z(x),
\]
where the free matrix $Z(x)$ is any matrix in
\[
\mathcal{Z}(x)=\{Z\in\mathbb{R}^{n\times n} : Zx=0\}.
\]
Given the mapping 
 \[
   \Phi:\mathcal{Z}(x) \to \mathbb{R},\quad \Phi(Z)=E(x;A_0+Z),
   \]
suppose that $\Phi(0_{n \times n}) \le 0$ and there exists $Z_1 \in\mathcal{Z}(x)$ such that $\Phi(Z_1)>0$.
\begin{comment}
    

Assume that further that the mapping
   \[
   \Phi:\mathcal{Z}(x) \to \mathbb{R},\quad \Phi(Z)=E(x;A_0+Z),
   \]
   depends nontrivally on $Z$.
 %  is continuously differentiable and there exists some $Z_0\in\mathcal{Z}(x)$ such that the derivative $D\Phi(Z_0):\mathcal{Z}(x) \to \mathbb{R}$ is surjective.
\end{comment}
   Then there exists a choice of $Z^*\in\mathcal{Z}(x)$ such that
\[
E(x;A_0(x)+Z^*)=0.
\]
%That is, for the given $x$ one may choose a semilinear representation $A(x)=A_0(x)+Z^*(x)$ (with $Z^*(x)x=0$) such that the SDRE residual vanishes.
   
\end{proposition}

\begin{proof}
    Since every semilinear form $A(x)$ satisfying $A(x)x=f(x)$ can be written as
   \[
   A(x)=A_0(x)+Z(x),
   \]
   with $Z(x)x=0$, the freedom in choosing $A(x)$ is equivalent to the freedom in choosing $Z(x)$ in the vector space
   \[
   \mathcal{Z}(x)=\{Z\in\mathbb{R}^{n\times n}: Zx=0\}.
   \]
   Note that if $x\neq 0$, then $\mathcal{Z}(x)$ is a linear subspace of dimension $n(n-1)$. With the above parameterization, the residual in the HJB equation can be viewed as a function
   \[
   \Phi:\mathcal{Z}(x)\to\mathbb{R},\quad \Phi(Z)=E(x;A_0(x)+Z).
   \]
  This mapping is continuously differentiable, as it results from the composition of smooth functions. Indeed, the residual $E(x;A_0(x)+Z)$ involves smooth operations and the SDRE solution $P(x;A_0(x)+Z)$ depends smoothly on its parameters (see, $e.g.$, \cite{ran1988parameter,sontag2013mathematical}).
  In the case $\Phi(0_{n \times n})=0$, the result follows by setting $Z^*=0_{n \times n}$. Otherwise, assume \(\Phi(0_{n \times n}) < 0\). By hypothesis, there exists \( Z_1 \in \mathcal{Z}(x) \) with \(\Phi(Z_1) > 0\).  
   Define the function \( f : [0,1] \to \mathbb{R} \) by
   $ 
   f(t) = \Phi(t Z_1).
   $
   Since scalar multiplication in $\mathcal{Z}(x)$ is continuous and \(\Phi\) is continuous, the composition \(f\) is continuous on \([0,1]\).  
   Notice that:
   \[
   f(0) = \Phi(0_{n \times n}) < 0 \quad \text{and} \quad f(1) = \Phi(Z_1) > 0.
   \]
   By the intermediate value theorem, there exists a \( t_0 \in (0,1) \) such that
   $
   f(t_0) = 0.
   $
   Setting \( Z^* = t_0 Z_1 \), which belongs to \(\mathcal{Z}(x)\) because \(\mathcal{Z}(x)\) is a linear subspace and is closed under scalar multiplication, we obtain
   \[
   \Phi(Z^*) = \Phi(t_0 Z_1) = 0,
   \]
   completing the proof.
    
   \end{proof}

    
%    Since $\Phi$ maps the (high--dimensional) linear space $\mathcal{Z}(x)$ to $\mathbb{R}$, the equation
%   $
%   \Phi(Z)=0
%   $
%   represents a single scalar equation in a space of dimension $n(n-1)$. Moreover, because $\Phi(Z)$ exhibits nontrivial dependence on $Z$, there exists at least one direction $\delta Z \in \mathcal{Z}(x) $ along which  $\Phi(Z)$ changes. Consequently, there exists some $Z_0\in \mathcal{Z}(x)$ such that differential $D\Phi(Z_0)$ is surjective.  By the implicit function theorem, the equation $\Phi(Z)=0$ admits a solution in a neighborhood of $Z_0$.   Therefore, there exists $Z^*\in\mathcal{Z}(x)$ such that
% $$
%   \Phi(Z^*)=E(x;A_0(x)+Z^*)=0.
%$$
%  \end{comment}


The aforementioned result is not merely theoretical, but also provides a systematic approach for constructing the optimal semilinear form. Given a reference semilinear form $A_0(x)$ and an alternative semilinear form $A_0(x)+Z(x)$, with $Z(x) \in \mathcal{Z}(x)$ and the corresponding residual $E(x)$ exhibits a sign change, an optimal semilinear form can be obtained by considering the continuous path $A_0(x)+tZ(x)$ for $t \in [0,1]$ and identifying the point where the residual vanishes.
For instance, fixing two pairs of indices $(i_1,j_1)$ and $(i_1,j_2)$, one may consider the following class of matrices in the set $\mathcal{Z}(x)$
\begin{equation}
[Z(x)]_{i,j}= \begin{cases}  x_{j_2} & (i,j)=(i_1,j_1),  \\
- x_{j_1} & (i,j)=(i_1,j_2), \\
0 & otherwise .
\end{cases}
\label{eq:Zx}
\end{equation}
It is straightforward to check that $Z(x)x = 0$ for all $x \in \mathbb{R}^d$. 
%The verification of the existence of these semilinear forms is difficult to state in general, one may consider different (possibly random) perturbation of the base semilinear form and verify directly the assumption.
\first{A theoretical characterization of when such semilinear forms exist is generally difficult. However, the construction becomes possible whenever the admissible class $\mathcal{Z}(x)$ contains at least one perturbation $Z(x)$ for which the associated residual, evaluated along the path $A_0(x) + t Z(x)$, changes sign. In practice, this can be assessed by introducing several (possibly random) perturbations of the baseline semilinear form and verifying whether such a sign change occurs.
}
%This idea has been already implemented 
%; it establishes the existence of an optimal semilinear form that ensures the residual $E(x)$ vanishes, but it does not provide a constructive method for deriving such a semilinear form. 
If identifying the sign-changing semilinear form proves challenging, the problem can be reformulated from an optimization perspective. Specifically, for a fixed $x \in \mathbb{R}^d$ and a base semilinear form $A_0(x)$, one can seek to minimize the residual norm by solving the optimization problem:
\begin{equation}
\min_{Z \in \mathcal{Z}(x)} | E(x;A_0+Z)|^2.
\label{min_E}
\end{equation}
One approach to tackling this problem is through a parametrization of the vector space $\mathcal{Z}(x)$. In \cite{jones2020} the authors attempt to solve the minimization problem \eqref{min_E} by considering a first-order Taylor expansion of the parametrization for $\mathcal{Z}(x)$, while in \cite{dolgov2022optimizing}, the authors restrict the minimization to a class of $\mathcal{Z}(x)$ formed using the class of linear perturbation described by \eqref{eq:Zx}. In both cases, the proposed numerical experiments demonstrate that the technique can generate a more stabilizing feedback. However, these methodologies are affected by dimensionality constraints and are primarily suitable for low-dimensional systems. \second{
\begin{remark}
Pursuing optimal or near-optimal semilinear forms remains possible in high dimension by exploiting additional structure.  
One option is to restrict $Z(x)$ to sparse, block-diagonal, or low-rank, whose complexity scales linearly rather than quadratically with the state dimension; see, e.g., \cite{benner2015survey}.  
Another strategy relies on projection-based model reduction: one computes an optimal semilinear form in a reduced-order model and lifts it back to the full space, following standard reduced-order optimal control techniques \cite{hinze2005podcontrol}.   

Alternatively, randomized and sampling-based approaches can be employed: by probing a small number of random directions in $\mathcal{Z}(x)$, one can efficiently detect sign changes of the residual $E(x;A_0+tZ)$ even in very large dimensions, leveraging ideas from high-dimensional randomized numerical linear algebra \cite{halko2011finding}.  
Finally, data-driven surrogate models may approximate the mapping $Z \mapsto E(x;A_0+Z)$ at fixed $x$, reducing the high-dimensional optimization to a learned low-complexity problem. Such approaches have been successful in the approximation of high-dimensional HJB-type equations \cite{onken2022neural,dolgov2023data}.
\end{remark}
}

\begin{example}[Allen-Cahn equation]
   We consider the following nonlinear Allen-Cahn PDE with homogeneous Neumann boundary conditions:
\begin{equation*}
\left\{ \begin{array}{l}
\partial_t y(t,x) = \sigma \partial_{xx} y(t,x) + y(t,x) (1-y(t,x)^2) + u(t,x),  \\
 y(0,x)=y_0(x),
\end{array} \right.
\end{equation*}
where $x \in [0,1]$ and $t \in (0,+\infty)$. The associated cost functional is given by
$$
J(u,y_0) = \int_0^{\infty}  \int_0^1 (|y(t,x)|^2 +  \tilde{\gamma}\|u(t,x)|^2) dx \, dt \,.
$$
By discretizing the PDE using finite difference schemes with $d$ grid points, we obtain the following system of ODEs:
\begin{align*}
\dot{y}(s)=A_0(y) y(s) + u(s),
\end{align*}
where
\begin{align*}
A_0(y) = \sigma A^{\Delta} + I_d -  diag(y \odot y),\quad y \in \mathbb{R}^d, 
\end{align*}
with $\odot$ denoting the Hadamard \first{(element-wise)} product, $I_d \in \mathbb{R}^{d \times d}$ being the identity matrix and $A^{\Delta}$ arising from the discretization of the Laplace operator with homogeneous Neumann boundary conditions. We set $\sigma = 10^{-q}$,  $d=100$, $\tilde{\gamma} = 0.1$ and $y_0(x) = \cos(\pi x)$. Our objective is to determine an optimal semilinear form $A(x)$ such that the residual $E(y_0;A(x))$ vanishes. To this end, we introduce a linear perturbation as defined in \eqref{eq:Zx}, specifically choosing $i_1=j_1=1$ and $j_2=2$, meaning that we perturb the first two entries of the first row of the semilinear form. We consider the parameterized form $A(x;\alpha) = A_0(x)+ \alpha Z(x)$, with $\alpha \in \mathbb{R}$. Left panel of Figure \ref{pics:AC} illustrates the behavior of the residual $E(y_0;\alpha)$ as the parameter $\alpha$ varies over the interval $[-10,10]$. The results exhibit a quadratic-like behavior, with at least one zero in the interval $(9,10)$, where a change of sign is observed. The root can be computed using the MATLAB function \texttt{fzero}, yielding $\alpha^* \approx 9.23 $, at which the residual attains the value $5.5 \cdot 10^{-15}$.
A similar result is obtained when considering the initial condition $y_0(x) = \sin (\pi x)$, as shown in the right panel of Figure \ref{pics:AC}, where the zero occurs at $\alpha^* \approx 5.74 $.

\begin{figure}[H]	
\begin{center}    
	\includegraphics[width=0.49\textwidth]{pics/AC/E_y0_AC.eps}	
    	\includegraphics[width=0.49\textwidth]{pics/AC/E_y0_AC2.eps}	
	
\caption{Allen-Cahn example. Variation of the residual as a function of the parameter $\alpha$ in the semilinear form, demonstrating the existence of at least one zero for $y_0(x) = \cos (\pi x)$ (left) and $y_0(x) = \sin (\pi x)$ (right).}
\label{pics:AC}
\end{center}


\end{figure}
\end{example}



\section{Numerical Methods for the Approximation of the SDRE}

In this section, we examine the numerical approximation of the SDRE methodology, highlighting that its implementation necessitates the repeated solution of the CAREs throughout the time integration process. This study investigates and compares two distinct numerical approaches for solving the sequence of CAREs, both of which adhere to Algorithm \ref{alg:sdre}.
For simplicity, we focus on the case $B(x)=B$. \first{We work under the assumptions of Theorem \ref{thm:stable}: 
$f(x)$ and $\frac{\partial f(x)}{\partial x_j} \ (j=1,\dots,n) $ are continuous near the origin,
$A(x)$ is continuous, and the pairs $ (A(x),B)$ and $(A(x),Q^{1/2})$  are, respectively, stabilizable and detectable
in a nonempty neighborhood of the origin.
}
%The first approach follows Algorithm \ref{alg:sdre}, wherein the Riccati equation is solved at each fixed state during the integration process. 
The first approach, known as the \emph{offline–online} method, employs a first-order approximation of the solution and precomputes specific components to mitigate computational costs during real-time execution. The second approach employs an iterative scheme, referred to as the \emph{Newton–Kleinman} method, in which the Riccati solution from the previous time step is used as an initial guess for the subsequent iteration. These two techniques are evaluated in terms of computational efficiency and accuracy through a numerical experiment involving the optimal control of a nonlinear reaction-diffusion PDE.
%This can be accomplished using either direct solution methods or iterative techniques. Here, we focus on the latter, as the solution from the previous time step can be utilized as an initial guess for the subsequent step. The second approach, referred to as the \emph{offline–online} method, involves precomputing specific components of the solution to mitigate computational costs during real-time implementation. Finally, we propose a \emph{hybrid} approach that combines the advantages of the offline–online framework with the computational efficiency of iterative schemes, aiming to improve both accuracy and efficiency in the numerical resolution of the SDRE methodology.



\subsection{Offline–Online Approach}

Suppose that the semilinear form admits the decomposition
\begin{equation}
  A(x) = A_0 + \sum_{j=1}^{r} f_j(x) A_j = A_0 + \tilde{A}(x),
  \label{eq:A_decomp}
\end{equation}
where $A_0,A_j \in \mathbb{R}^{d \times d}$ and $f_j: \mathbb{R}^d \rightarrow \mathbb{R}$ for $j \in \{1,\dots, r\}$. This formulation is commonly encountered in the optimal control of PDEs, where the system dynamics include both a linear operator (e.g., diffusion or convection) and a nonlinear term, such as a reaction component.

Within this framework, the solution to the Riccati equation can be approximated as
\begin{equation}
  P(x) \approx P_0 + \sum_{j=1}^{r} P_j f_j(x) = P_0 + W(x),
  \label{eq:pi_approx}
\end{equation}
where:
\begin{itemize}
  \item \(P_0\) is computed by solving the CARE for the linearized system corresponding to \(A_0\):
    \begin{equation}
    A_0^\top P_0 + P_0 A_0 - P_0 S P_0 + Q = 0,
    \label{P0}
        \end{equation}
    with $
    S = B R^{-1}B^\top.
   $
  \item Each \(P_j\) is obtained by solving an associated Lyapunov equation
    \[
    P_j C_0 + C_0^\top P_j + Q_j = 0, \quad j=1,\ldots,r,
    \]
    with \(C_0 = A_0 - SP_0\) and \(Q_j = P_0 A_j + A_j^\top P_0\).
\end{itemize}




This methodology remains computationally feasible when the number of scalar nonlinear terms, $r$ is small. However, in many practical scenarios $r$ is comparable to the system dimension $d$, which can be arbitrarily large, making the direct implementation impractical due to the excessive computational burden.
When both $r$ and $d$ are large, an \emph{offline–online} strategy can be adopted to mitigate computational complexity. During the offline phase, the matrix $P_0$ is precomputed by solving \eqref{P0}. In the online phase, for a given state \(x\), the term $W(x)$ is computed by solving a single Lyapunov equation:
\begin{equation}
  W(x) C_0 + C_0^\top W(x) + Q(x) = 0,
  \label{eq:W_lyap}
\end{equation}

where 
$$
Q(x) = \sum_{j=1}^{r} Q_j f_j(x) = P_0 \tilde{A}(x) + \tilde{A}(x)^\top P_0.
$$
The resulting feedback law is then given by
\begin{equation}
  u_{O}(x) = -R^{-1}B^\top\Bigl( P_0 + W(x) \Bigr)x.
  \label{eq:feedback_offline_online}
\end{equation}




\begin{algorithm}
\caption{Offline-online SDRE}
\label{alg:off_onl}{
\begin{algorithmic}[1]
\Statex{\textbf{Input:} $A(\cdot), B, Q,R$, initial conditions $ y_0$, time discretization $\{t_i\}_{i=0}^{N_T}$ }
\Statex{\textbf{Output:} Controlled trajectory $\{y_O(t_i)\}_{i=0}^{N_T}$}
\Statex{}
\Statex \textbf{Offline phase}
\State Compute the solution $P_0$ from \eqref{P0}
\Statex{}
\Statex \textbf{Online phase}
\State{$y_O(t_0):= y_0$}
\For{$n = 0, \dots, N_T-1$}
\State{Compute $W(y_O(t_{n}))$ from \eqref{eq:W_lyap} }
\State{Build the control $u_O(y_O(t_{n}))$ from \eqref{eq:feedback_offline_online}}
\State{Integrate the system \eqref{semilinear} to obtain $y_O(t_{n+1})$}       
\EndFor 
\end{algorithmic}}
\end{algorithm}

The method is sketched in Algorithm \ref{alg:off_onl}.
This approach is attractive in real–time applications because it decouples the heavy computation (solved offline) from the online computation that is reduced to the resolution of one Lyapunov equation per time step. However, the stability of the closed-loop matrix $C(x) = A(x)-S\Bigl(P_0+W(x)\Bigr)$ is not inherently guaranteed. The following proposition provides a rigorous framework under which this approach ensures the stability of the closed-loop matrix.


\begin{proposition}\label{prop:stability}
Given the semilinear form \eqref{eq:A_decomp} and $P_0$ the solution of the CARE \eqref{P0}. Let $W(x)$ be the unique solution of the Lyapunov equation \eqref{eq:W_lyap} and let $C_0$ be diagonalizable with decomposition $C_0 = V \Lambda V^{-1}$.

Then if
\begin{equation}\label{eq:bound2}
\left\|\tilde{A}(x)\right\| \left(1+\|S\|\frac{M^2}{\alpha}\left\|P_0\right\|\right) < \alpha,
\end{equation}
where $\alpha = \min_{\lambda \in \sigma(C_0)} |\operatorname{Re}(\lambda)|$, with $\operatorname{Re}(\cdot)$ denoting the real part operator, and $M = cond(V) = \|V \| \| V^{-1} \|$,
the closed-loop matrix $C(x) = A(x)-S\Bigl(P_0+W(x)\Bigr)$ is Hurwitz \first{(i.e.\ with eigenvalues with negative real part)}.
\end{proposition}


\begin{proof}

The solution $W(x)$ of the Lyapunov equation \eqref{eq:W_lyap} can be expressed as
$$
W(x)=- \int_{0}^{\infty} e^{C_0^\top t} Q(x) e^{C_0 t}\,dt.
$$

Since $P_0$ is the unique solution of the CARE \eqref{P0}, the corresponding closed-loop matrix $C_0$ is Hurwitz. Consequently,
$$
\|e^{C_0 t}\| = \|V e^{\Lambda t } V^{-1}\| \le M e^{-\alpha t},
$$
with $\alpha = \min_{\lambda \in \sigma(C_0)} |\operatorname{Re}(\lambda)|$.
Applying this bound to the integral representation of $W(x)$, we obtain
$$
\|W(x)\| \le M^2 \| Q(x) \| \int_0^{+\infty}  e^{-2 \alpha t} \,dt = \frac{M^2}{2\alpha}  \| Q(x) \|. 
$$
Furthermore, we note that 
$$
\|Q(x) \| \le 2 \|P_0\| \|\tilde{A}(x)\|.
$$

Noting that $C(x) = C_0 + C_1(x) = C_0+ \tilde{A}(x)-S W(x)$,
a sufficient condition for the closed-loop matrix $C(x)$ to be stable is that the norm of $C_1(x)$ remains strictly less than the stability margin $\alpha$ of $C_0$. By applying the derived bound on $\|W(x)\|$ and $\|Q(x)\|$, this requirement leads to the computable stability criterion given by \eqref{eq:bound2}.
 
\end{proof}
%\textcolor{red}{a proof that if $A(x)-A_0$ is smalll, then $P(x)$ is a stabilizing solution.}

The result presented above establishes a condition under which the closed-loop matrix $C(x)$ remains Hurwitz, thereby ensuring the synthesis of a stabilizing feedback law. However, if the norm of the nonlinear perturbation $\|\tilde{A}(x)\|$ is excessively large, the stability of the closed-loop matrix may be compromised. In such cases, the resulting controlled trajectory may fail to converge to the desired equilibrium. In the following section, we introduce an alternative approach based on an iterative scheme, which, although possibly requiring the solution of additional Lyapunov equations, ensures the stability of the closed-loop matrix under less restrictive conditions.


\subsection{Newton–Kleinman Method}
An alternative approach relies on the resolution of the full SDRE \eqref{sdre} at each state of the dynamical system, thereby avoiding the approximation error introduced by the decomposition in \eqref{eq:pi_approx}. Since the SDRE must be solved at multiple time steps, an efficient strategy involves employing an iterative scheme, such as the Newton–Kleinman (NK) method.  At a given time step $t_n$ for a fixed state $x$, the NK method refines an initial guess \(P^{(0)}(x)\) iteratively, updating it according to the recurrence relation
\begin{equation}
  P^{(k+1)}(x) = P^{(k)}(x) + \Delta P(x),
  \label{eq:NK_update}
\end{equation}
where the correction \(\Delta P(x)\) is computed by solving the Lyapunov equation
\begin{equation}
  (A(x)-SP^{(k)}(x) )^\top \Delta P(x) + \Delta P(x) (A(x)-SP^{(k)}(x) )  = -\mathcal{R}(P^{(k)},x),
  \label{eq:NK_lyap}
\end{equation}
with the residual
\[
\mathcal{R}(P^{(k)},x) = A^\top(x) P^{(k)}(x) + P^{(k)}(x) A(x) - P^{(k)}(x) S P^{(k)}(x) + Q.
\]
The procedure is repeated until the residual $\mathcal{R}(P^{(\tilde{k})},x)$ falls below a prescribed threshold or the maximum number of iterations is reached. Subsequently, we assign $P_{n}(x):=P^{(\tilde{k})}(x)$.
In the context of time–varying nonlinear systems, one can \emph{warm start} the NK iteration by using the Riccati solution from the previous time step as the initial guess \(P^{(0)}\). The technique is outlined in Algorithm \ref{alg:CNK}.
This strategy, which leverages the temporal continuity of the Riccati solution to enhance computational efficiency and convergence, is referred to as the \emph{cascade Newton–Kleinman} (\texttt{C-NK}) method. In \cite{saluzzi2025dynamical}, the authors demonstrate that if the time step is sufficiently small, bounded by a computable threshold, then the solution from the previous iteration $P(y_S(t_{n-1}))$ stabilizes the pair $(A(y_S(t_n)),B)$, $i.e.$, the closed-loop matrix $A(y_S(t_n))-S P(y_S(t_{n-1}))$ is stable, ensuring the convergence of the NK scheme (see \cite{kleinman1968} for a precise formulation).

The warm start strategy can be applied only from the second time step onward, as the initial Riccati solution at the first time step cannot be obtained from a previous iteration. Consequently, the first instance of the SDRE \eqref{sdre} is solved using a direct method, such as the MATLAB function \texttt{icare}. This step is analogous to the first phase of the offline-online method, where the initial Riccati equation \eqref{P0} must be solved. \first{The method does not require the input matrix $B$ to be constant; it can be applied as well when $B$ is state dependent, namely $B = B(x)$.
}

The key distinction between these two methods lies in their computational requirements: while the offline-online approach necessitates solving a single Lyapunov equation per time step, the cascade Newton–Kleinman scheme may require multiple iterations to achieve convergence. The availability of a suitable initial guess in the \texttt{C-NK} method not only accelerates convergence but also ensures a stabilizing initial condition, provided that the time step is sufficiently small.
In the following section, we present a numerical experiment demonstrating the effectiveness of this technique.


\begin{algorithm}
\caption{\texttt{C-NK} SDRE}
\label{alg:CNK}{
\begin{algorithmic}[1]
\Statex{\textbf{Input:} $A(\cdot), B, Q,R$, initial conditions $ y_0$, time discretization $\{t_i\}_{i=0}^{N_T}$}
\Statex{\textbf{Output:} Controlled trajectory $\{y_S(t_i)\}_{i=0}^{N_T}$}
\State{$y_S(t_0):= y_0$}
\State{Compute $P_0 = P(y_0)$ from \eqref{sdre}}
\State{Compute the control $u_S(y_0)$ from \eqref{control_sdre}}
\State{Integrate the system \eqref{semilinear} to obtain $y_S(t_{1})$}
\For{$n = 1, \dots, N_T-1$}
\State{Compute $P_n$ from \eqref{eq:NK_update}-\eqref{eq:NK_lyap} with guess $P_{n-1}$ }
\State{Build the control $u_S(y(t_{n}))$ from \eqref{control_sdre}}
\State{Integrate the system \eqref{semilinear} to obtain $y_S(t_{n+1})$}       
\EndFor 
\end{algorithmic}}
\end{algorithm}
%Although warm–starting often reduces the number of iterations required for convergence, the NK method may still necessitate multiple Lyapunov solves per time step. To address this limitation, in the subsequent section, we exploit the advantages of the \emph{Offline–Online} approach to construct a more accurate initial guess for the NK scheme. By leveraging precomputed components in the offline phase, this method aims to further reduce the number of iterations required in the online phase, thereby enhancing computational efficiency without compromising accuracy.

\begin{comment}

\subsection{Hybrid Method}
The hybrid method is designed to leverage the fast online prediction of the offline–online approach and the accuracy improvements from a few NK iterations. The idea is as follows:
\begin{enumerate}
  \item \textbf{Predictor Step:} Compute the approximate Riccati solution using the offline–online method:
    \begin{equation}
      P_{\text{pred}} = P_0 + W(x),
      \label{eq:predicted_pi}
    \end{equation}
    and obtain the corresponding predicted gain
    \begin{equation}
      K_{\text{pred}}(x) = -R^{-1}B^\top P_{\text{pred}}.
      \label{eq:predicted_gain}
    \end{equation}
  \item \textbf{Corrector Step:} Use \(P_{\text{pred}}\) as the initial guess for a small number of NK iterations:
    \[
    P^{(0)} = P_{\text{pred}},
    \]
    and update according to \eqref{eq:NK_update} and \eqref{eq:NK_lyap} until either the residual \(\|\mathcal{R}(P^{(k)})\|\) is below a specified tolerance or a maximum number \(N_{\text{iter}}\) of iterations is reached.
\end{enumerate}
The final feedback gain is then computed as
\begin{equation}
  K(x) = -R^{-1}B^\top P^{(N_{\text{iter}})}.
  \label{eq:hybrid_gain}
\end{equation}
This hybrid strategy capitalizes on the efficiency of the offline–online approach while allowing for corrective NK iterations to improve the solution accuracy when necessary. In many practical scenarios, the predictor \(P_{\text{pred}}\) is already sufficiently accurate, so that only one or two NK iterations are needed—thus maintaining a low computational cost per time step.

\end{comment}





\subsection{Numerical Example}

Throughout this discussion, we denote by $\chi_{\omega}(\cdot)$ the indicator function over the subdomain $\omega \subset \mathbb{R}$. The interval $\omega_c$ specifies the region where the control input is applied, while $\omega_o$ designates the observation domain where the system dynamics are monitored.
We consider the nonlinear Zeldovich-type equation with homogeneous Neumann boundary conditions:
\begin{equation*}
\left\{ \begin{array}{l}
\partial_t y(t,x) = \sigma \partial_{xx} y(t,x) + \nu y(t,x)+ \mu y(t,x)^2 (1-y(t,x)) + u(t,x)\chi_{\omega_c}(x),  \\
 y(0,x)=y_0(x),
\end{array} \right.
\label{AC}
\end{equation*}
with $x \in [0,1]$ and $t \in (0,+\infty)$. The corresponding cost functional is given by
$$
 J(u,y_0) = \int_0^{+\infty} \int_{\omega_o} | y(t,x)|^2 \, dx + \tilde{\gamma} \int_{\omega_c} | u(t,x) |^2 \, dx \, dt.
$$
By discretizing the PDE using a finite difference scheme with $d$ grid points, we obtain the following system of ODEs 
\begin{align*}
\dot{y}(s)=A(y) y(s) + u(s),
\end{align*}
with
$$
A(y) = A_0 +  \mu \, diag(y-y \odot y),\quad y \in \mathbb{R}^d,
$$
$$
A_0 = \sigma A^{\Delta} + \nu I_d.
$$
Here, $\odot$ denotes the Hadamard product, $I_d \in \mathbb{R}^{d \times d}$ is the identity matrix and $A^{\Delta}$ is the discrete Laplacian matrix arising from the finite difference approximation with Neumann boundary conditions. We set the discretization dimension to $d =100$ and initialize the system with $y_0(x) = \cos (\pi x)$. The dynamical system is discretized using a semi-implicit first-order integration scheme, where the diffusion term is treated implicitly, while the reaction and control terms are handled explicitly. A fixed time step of $\Delta t = 0.02$ is employed, with the final simulation time set to $T=4$, leading to a total of $200$ time steps. We conduct a comparative analysis of the offline-online method (Algorithm \ref{alg:off_onl}), the \texttt{C-NK} method (Algorithm \ref{alg:CNK}), and the "\texttt{icare} approach", which is based on Algorithm \ref{alg:sdre} and utilizes the MATLAB function \texttt{icare} to solve the Riccati equations. In the \texttt{C-NK} method, the stopping criterion for the residual is set to $10^{-5}$. Two different parameter configurations are considered:
\begin{enumerate}
    \item $\sigma = 0.2$, $\tilde{\gamma} = 0.01$, $\nu = 0.5$ where the control is applied in $\omega_c= [0.2,0.5] $ and the observation is restricted to $ \omega_o = [0.5,0.7]$, leading to a partial-domain control and monitoring setup.
    \item $\sigma = 10^{-2}$, $\tilde{\gamma} = 0.1$, $\nu = 0.5$ with full-domain control and observation regions $\omega_c = \omega_o = [0,1]$, meaning that both the control and observation mechanisms cover the entire spatial domain.
\end{enumerate}
These two scenarios exhibit distinct behaviors in the decay of the singular values of the Riccati solution. Fixing $\mu=1$, we solve the corresponding SDRE at the initial condition $y_0$. The decays of the singular values for both test cases are illustrated in the left panel of Figure \ref{pics:sing_value}. In the first case, the singular values exhibit a rapid decay, suggesting a low-rank structure in the solution. In the second case, the singular values decay slowly, indicating that the Riccati solution is full-rank. For high-dimensional problems, these two cases require different numerical treatments. For low-rank solutions, there exists a substantial body of literature on efficient methods for solving low-rank Lyapunov and Riccati equations. For further details, we refer the reader to \cite{simoncini2014two,simoncini2016computational,benner2008numerical}. Conversely, in the full-rank setting, the sparsity and potential banded structure of the solution can be effectively leveraged, as illustrated in the right panel of Figure \ref{pics:sing_value}. For further details, the reader is referred to \cite{palitta2018,massei2024data,haber2018sparsity} and the references therein. A comprehensive analysis and implementation of these advanced techniques fall beyond the scope of this study. Given the problem dimension considered, the Lyapunov equations are solved using the MATLAB function \texttt{lyap}, while the initial Riccati equation is computed employing the MATLAB command \texttt{icare}.




\begin{figure}[H]	
\begin{center}    
	\includegraphics[width=0.49\textwidth]{pics/AC/sing_values_full_low.eps}	
	\includegraphics[width=0.49\textwidth]{pics/AC/Py0.eps}
\caption{Zeldovich example. \emph{Left}: Decay of the singular values of the Riccati solution $P(y_0)$ for the two different cases. \emph{Right}: Off diagonal decay for the solution $P(y_0)$ in the second case.}
\label{pics:sing_value}
\end{center}


\end{figure}

\vspace{0.5cm}

{\bf Case 1}

\vspace{0.2cm}

We first consider the scenario where the control and observation domains are given by $\omega_c= [0.2,0.5] $ and $ \omega_o = [0.5,0.7]$, respectively, with parameters $\sigma = 0.2$, $\tilde{\gamma} = 0.01$ and $\nu = 0.5$. The reaction parameter $\mu$ varies within the set $\{1,2\}$. Table \ref{table_AC_low_rank} presents the CPU time and the total cost obtained using different numerical approaches for the considered values of $\mu$. The results indicate that the most efficient method, achieving the lowest total cost, is the \texttt{C-NK} approach. While the SDRE method based on \texttt{icare} yields a total cost comparable to that of \texttt{C-NK}, being approximately 60 times slower. Conversely, the offline-online approach requires up to four times the CPU time and results in a higher total cost.  The superior computational efficiency of \texttt{C-NK} can be attributed to the fact that, for certain time steps, the previously computed solution exhibits a low residual. As a result, the resolution of a Lyapunov equation is not required at every step, reducing computational overhead.
In the case $\mu=2$, the control law generated by the offline-online method fails to stabilize the system, as illustrated in the right panel of Figure \ref{pics:AC_cost_low_rank}.
Furthermore, Figure \ref{pics:AC_cost_low_rank} demonstrates that the \texttt{C-NK} and \texttt{icare}-based methods exhibit similar performance trends.


\begin{table}[hbht]
\centering
\begin{tabular}{c|cc|cc|cc}    

        & \multicolumn{2}{c|}{offline-online}   & \multicolumn{2}{c|}{\texttt{C-NK}} & \multicolumn{2}{c}{\texttt{icare}} \\
$\mu$   & CPU time & Total cost & CPU time & Total cost & CPU time & Total cost 
\\\hline
1    & 0.98  &   3.10e-01 & 0.27   &  3.08e-01  &  15.7 & 3.08e-01\\

2 & 1.02  &   1.63e+03  & 0.25   &  8.56e-01   &  15.5 & 8.56e-01  \\
 \end{tabular}
  \caption{Zeldovich example (Case 1). Comparison of the performance in the computation of the controlled trajectory via the offline-online approach, \texttt{C-NK} method and via the \texttt{icare}-based SDRE.}
 \label{table_AC_low_rank}
\end{table}


\begin{figure}[H]	
\begin{center}    
	\includegraphics[width=0.49\textwidth]{pics/AC/cost_mu1_low_rank.eps}	
	\includegraphics[width=0.49\textwidth]{pics/AC/cost_mu2_low_rank.eps}	
\caption{Zeldovich example (Case 1). Running cost for the controlled trajectories with the different techniques with $\mu =1$ (left) and $\mu=2$ (right).}
\label{pics:AC_cost_low_rank}
\end{center}


\end{figure}


\vspace{0.5cm}

{\bf Case 2}

\vspace{0.2cm}

In this second case, we consider the fully controlled and observed setting, where $\omega_c = \omega_o = [0,1]$. The parameters are set to $\sigma = 10^{-2}$, $\tilde{\gamma} = 0.1$ and $\nu = 0.5$, while the reaction coefficient $\mu$ varies within the set $\{1,2\}$. The results obtained for the three numerical techniques across different values of $\mu$ are summarized in Table \ref{table_AC_full_rank} and illustrated in Figure \ref{pics:AC_cost_full_rank}. For the choice $\mu = 2$, the final simulation time is reduced to 0.6, due to the fast divergence of the offline-online method.
The observed trends are consistent with those in the previous case. Once again, the \texttt{C-NK} method demonstrates the best performance in terms of both computational efficiency and total cost. The \texttt{icare}-based SDRE approach achieves the same total cost as \texttt{C-NK}, but its computational time is more than 40 times higher for the first choice of $\mu$, further emphasizing the efficiency of \texttt{C-NK}.
Regarding the offline-online method, for $\mu =1$ the total cost is comparable to that of the other two techniques. As in the previous case, increasing the reaction coefficient to $\mu = 2$ results in the failure of the offline-online approach to stabilize the system.


\begin{table}[hbht]
\centering
\begin{tabular}{c|cc|cc|cc}    

        & \multicolumn{2}{c|}{offline-online}   & \multicolumn{2}{c|}{\texttt{C-NK}} & \multicolumn{2}{c}{\texttt{icare}} \\
$\mu$   & CPU time & Total cost & CPU time & Total cost & CPU time & Total cost 
\\\hline
1    & 0.53  &   2.49e-01 & 0.46   &  2.38e-01  &  21.2 & 2.38e-01\\

2 & 0.16  &   3.57e+04 & 0.30   &  2.78e-01   &  8.11 & 2.78e-01  \\
 \end{tabular}
  \caption{Zeldovich example (Case 2). Comparison of the performance in the computation of the controlled trajectory via the offline-online approach, \texttt{C-NK} method and via the \texttt{icare}-based SDRE.}
 \label{table_AC_full_rank}
\end{table}


\begin{figure}[H]	
\begin{center}    
	\includegraphics[width=0.49\textwidth]{pics/AC/cost_mu1_full_rank.eps}	
    \includegraphics[width=0.49\textwidth]{pics/AC/cost_mu2_full_rank.eps}	
\caption{Zeldovich example (Case 2). Running cost for the controlled trajectories with the different techniques with $\mu =1$ (left) and $\mu=2$ (right).}
\label{pics:AC_cost_full_rank}
\end{center}
\end{figure}

%To better study this behavior, we check the condition \eqref{eq:bound} given by Proposition \ref{prop:stability}. We note that in this case $C_0=A_0-S P_0$ is symmetric, since $A_0$ is symmetric and $S$ is the multiplication of the identity for a constant.

\section{Conclusions}

    
This work provides an analysis of the State-Dependent Riccati Equation method for nonlinear optimal control, focusing on its theoretical foundations, numerical approximations, and computational efficiency. We examined the relationship between SDRE and the HJB equation, deriving error estimates based on the residual introduced by the SDRE approximation. Our analysis highlighted the role of the semilinear decomposition in influencing the accuracy of the SDRE approach and proposed an optimal decomposition strategy to minimize approximation errors.

From a computational perspective, we compared two numerical methods for solving the SDRE: the offline–online approach and the Newton–Kleinman iterative method. Numerical experiments demonstrated that while the offline–online approach offers a reduced computational burden, it may fail to stabilize the system under certain conditions. In contrast, the C-NK method consistently provided more accurate and stable controlled solutions, making it a preferable choice in real-time control applications.

Future work could explore more efficient numerical methods for solving the sequence of Riccati equations in high-dimensional settings, leveraging low-rank approximations, sparsity-preserving methods, or data-driven techniques. Additionally, extending the SDRE framework to stochastic control scenarios would be a valuable direction for further research.

\vspace{1cm}

{\bf Acknowledgements.} The work has been partially supported by Italian Ministry of
Instruction, University and Research (MIUR) (PRIN Project 2022 2022238YY5, “Optimal control problems: analysis, approximation”) and INdAM-research group GNCS (CUP E53C24001950001,
“Metodi avanzati per problemi di Mean Field Games ed applicazioni”).

\vspace{5mm}

{\bf Data Availability.} Data will be made available on request.


\begin{comment}
    

\begin{appendices}
\section{Proof of Theorem \ref{thm:stable2}}
\label{appendixA}

Under the feedback control
$     \tilde{u}_{S}(x) $, the closed-loop dynamics becomes
\begin{equation}
    \dot{x} = A(x)x - B(x) R^{-1} B(x)^T P(x)x - \frac{1}{2} B(x) R^{-1} B(x)^T \varphi(x).
    \label{closed_loop}
\end{equation}
We decompose the state-dependent matrices as
$$
    A(x) = A_0 + \tilde{A}(x), \quad B(x) = B_0 + \tilde{B}(x), \quad P(x) = P_0 + \tilde{P}(x),
$$
where \( A_0 = A(0) \), \( B_0 = B(0) \), \( P_0 = P(0) \), and \( \tilde{A}(0) = \tilde{B}(0) = \tilde{P}(0) = 0_{d \times d} \). Also, by definition, \( \varphi(0) = 0 \).
We note that $x=0$ is an equilibrium for the closed-loop dynamics \eqref{closed_loop}.

Neglecting higher-order terms, the linearization of the closed-loop system is
$$
   A_{cl,0} \, x = (A_0 - B_0 R^{-1} B_0^T P_0) x,
$$
where $ A_{cl,0}$ is Hurwitz.

Define the perturbation term
$$
    h(x) = \Bigl[ \tilde{A}(x)x - \tilde{B}(x)R^{-1}B(x)^T P(x)x - B_0R^{-1}B_0^T\tilde{P}(x)x \Bigr] - \frac{1}{2} B(x)R^{-1}B(x)^T\varphi(x).
$$
Since \( \tilde{A}(x) \), \( \tilde{B}(x) \), and \( \tilde{P}(x) \) vanish at \( x=0 \) and are continuous, and because \( \varphi(x) \) is at least quadratic in \( x \), it follows that
\begin{equation}
    \lim_{x \to 0} \frac{\|h(x)\|}{\|x\|} = 0.
\end{equation}
Thus, the closed-loop dynamics can be written as
$$
    \dot{x} = A_{cl,0} x + h(x).
$$
Applying the variation-of-constants formula and Gronwall’s inequality, we obtain
$$
    \|x(t)\| \le G e^{-(\beta - \zeta_0 G)t} \|x_0\|,
$$
where \( \beta - \zeta_0 G > 0 \). This ensures exponential decay to the origin, proving local asymptotic stability.

Thus, the control $\tilde{u}_{S}(x)$ renders the closed-loop system locally asymptotically stable at the origin.

\end{appendices}
\end{comment}
\vspace{1cm}

\bibliographystyle{spmpsci}
\bibliography{sample,biblio}
\end{document}






